\documentclass[11pt,a4paper]{article}
\usepackage[margin=1.8cm]{geometry}
\usepackage{amsmath,amssymb,array,tabularx,fontspec,cancel,tikz,graphicx}
\setmainfont{FreeSerif}
\newfontfamily\urdufont[Script=Arabic]{FreeSerif}
\TeXXeTstate=1
\newcommand{\urdu}[1]{{\urdufont\beginR #1\endR}}
\newcommand{\lr}[1]{\beginL #1\endL}
\setlength{\parindent}{0pt}
\setlength{\parskip}{0pt}
\renewcommand{\arraystretch}{1.15}
\newcommand{\ut}[1]{\text{{\urdufont\beginR #1\endR}}}
% \uq{N}{Urdu text}: question with its printed number on the right
\newcommand{\uq}[2]{\par\medskip{\urdufont\raggedleft\beginR\textbf{\lr{#1.}}\ #2\endR\par}}
% continuation paragraph (no number)
\newcommand{\usc}[1]{\par\smallskip{\urdufont\raggedleft\beginR #1\endR\par}}
% centred math line
\newcommand{\mline}[1]{\par\smallskip\centerline{\lr{$\displaystyle #1$}}}
\newcommand{\tfig}[1]{\par\smallskip\centerline{{\TeXXeTstate=0 \begin{tikzpicture}[scale=0.8,font=\small]#1\end{tikzpicture}}}\par}
\newcommand{\pfig}[2]{\par\smallskip\centerline{\includegraphics[width=#2]{#1}}\par}
% data table (LTR table inside a centred line)
\newcommand{\dtab}[2]{\par\smallskip\centerline{\lr{\begin{tabular}{#1}\hline #2\\ \hline\end{tabular}}}\par\smallskip}
\newcommand{\nums}[1]{\par\smallskip\centerline{\lr{\parbox{0.9\linewidth}{#1}}}\par\smallskip}
\newcommand{\foodsurvey}{ایک سروے کیا گیا جس میں \lr{130} لوگوں سے ان کے پسندیدہ کھانوں کے بارے میں پوچھا گیا۔ سروے کے نتائج نے درج ذیل اعداد و شمار ظاہر کیے۔ \lr{40} لوگوں نے کہا کہ وہ نہاری کو پسند کرتے ہیں۔ \lr{65} لوگوں نے کہا کہ انہیں بریانی پسند ہے۔ \lr{50} لوگوں نے کہا کہ انہیں قورمہ پسند ہے۔ \lr{20} لوگوں نے کہا کہ انہیں نہاری اور بریانی پسند ہے۔ \lr{35} لوگوں نے کہا کہ انہیں بریانی اور قورمہ پسند ہے۔ \lr{27} لوگوں نے کہا کہ انہیں نہاری اور قورمہ پسند ہے۔ \lr{12} لوگوں نے کہا کہ انہیں تینوں کھانے نہاری، بریانی اور قورمہ پسند ہیں۔}
\newcommand{\mallemp}{ایک شاپنگ مال میں \lr{1} سے \lr{150} لیبل والے \lr{150} ملازمین ہیں جو یونیورسل سیٹ \lr{U} کو ظاہر کرتے ہیں۔ ملازمین درج ذیل زمروں میں آتے ہیں: سیٹ \lr{A: 40} ملازمین جن کی تنخواہ کی حد \lr{30} ہزار تا \lr{45} ہزار ہے۔ جس پر \lr{50} سے \lr{89} تک کا لیبل لگایا گیا ہے۔ سیٹ \lr{B: 50} ملازمین جن کی تنخواہ کی حد \lr{50} ہزار تا \lr{80} ہزار ہے۔ جس کا لیبل \lr{101} سے \lr{150} تک ہے۔ سیٹ \lr{C: 60} ملازمین جن کی تنخواہ کی حد \lr{100} ہزار تا \lr{150} ہزار ہے۔ \lr{1} سے \lr{49} اور \lr{90} سے \lr{100} تک کا لیبل لگا ہوا ہے۔}
\newcommand{\sportsch}{ایک سیکنڈری سکول میں \lr{125} طلبہ میں سے ہر ایک درج ذیل میں سے کم از کم ایک کھیل میں حصہ لیتا ہے۔ کرکٹ، فٹ بال یا ہاکی۔ \lr{60} طلبہ کرکٹ کھیلتے ہیں۔ \lr{70} طلبہ فٹ بال کھیلتے ہیں۔ \lr{40} طلبہ ہاکی کھیلتے ہیں۔ \lr{25} طلبہ کرکٹ اور فٹ بال دونوں کھیلتے ہیں۔ \lr{15} طلبہ فٹ بال اور ہاکی دونوں کھیلتے ہیں۔ \lr{10} طلبہ کرکٹ اور ہاکی دونوں کھیلتے ہیں۔}
\newcommand{\itvenn}{ملحقہ وین شکل پر غور کریں جس میں ایک آئی ٹی ادارے میں مختلف کورسز میں داخلہ لینے والے طلبہ کو ظاہر کیا گیا ہے۔ \lr{U=} آئی ٹی ادارے میں داخلہ لینے والے طلبہ، \lr{A=} اپلائیڈ روبوٹکس میں داخلہ لینے والے طلبہ، \lr{G=} گیم ڈویلپمنٹ میں داخلہ لینے والے طلبہ \lr{W=} ویب ڈیزائننگ میں داخلہ لینے والے طلبہ۔}
\newcommand{\shopsurvey}{ایک شاپنگ مال میں \lr{130} صارفین کا سروے کیا گیا۔ جس میں ان سے خریداری کی ترجیحات کے بارے میں پوچھا گیا۔ سروے کے نتائج نے درج ذیل اعداد و شمار ظاہر کیے۔ \lr{57} صارفین نے گارمنٹس کی خریداری کی۔ \lr{50} صارفین نے سنگھار کا سامان خریدا۔ \lr{46} صارفین نے الیکٹرونکس کی اشیا خریدی۔ \lr{31} صارفین نے گارمنٹس اور سنگھار دونوں کا سامان خریدا۔ \lr{25} صارفین نے گارمنٹس اور الیکٹرانکس دونوں کی خریداری کی۔ \lr{21} صارفین نے سنگھار اور الیکٹرونکس دونوں کی خریداری کی۔ \lr{12} صارفین نے تینوں مصنوعات یعنی گارمنٹس، سنگھار اور الیکٹرونکس کی خریداری کی۔}
\newcommand{\fzA}[2]{\uq{#1}{تجزی کریں: \lr{$#2$}}}
\newcommand{\fzB}[2]{\uq{#1}{درج ذیل جملے کی تجزی کریں۔ \lr{$#2$}}}
\newcommand{\abprop}{ذیل میں دیے گئے سیٹ \lr{A,B} اور \lr{C} کے لیے خصوصیات کو ثابت کریں: (الف) یونین کی خاصیت تلازم (ب) تقاطع کی خاصیت تلازم (ج) یونین کی تقاطع پر خاصیت تقسیمی (د) تقاطع کی یونین پر خاصیت تقسیمی}
\newcommand{\examtxt}{سالانہ امتحان میں شامل ہونے والے \lr{180} طلبہ میں سے \lr{120} نے ریاضی کا امتحان پاس کیا، \lr{90} نے سائنس کا امتحان پاس کیا۔ اور \lr{60} نے ریاضی اور سائنس دونوں ٹیسٹ پاس کیے۔}
\newcommand{\hcfdiv}[2]{\uq{#1}{درج ذیل جملے کا عاد اعظم بذریعہ تقسیم معلوم کریں:}\mline{#2}}
\newcommand{\hcffac}[2]{\uq{#1}{عاد اعظم بذریعہ تجزی معلوم کریں: \lr{$#2$}}}
\newcommand{\sqrdiv}[2]{\uq{#1}{درج ذیل جملے کا جذر المربع بذریعہ تقسیم معلوم کریں:}\mline{#2}}
\newcommand{\shade}[3]{\uq{#1}{درج ذیل یک درجی غیر مساواتوں کے حل کا خطہ سایہ دار کر کے ظاہر کریں: \lr{$#2,\ #3$}}}
\newcommand{\trig}[2]{\uq{#1}{درج ذیل تکونیاتی اکائی کو ثابت کریں:}\mline{#2}}
\newcommand{\solvetri}[2]{\uq{#1}{درج ذیل مثلث کو حل کریں جب کہ: \lr{$#2$}}}
\newcommand{\quarter}[2]{\uq{#1}{اگر \lr{$\sin\frac{\pi}{4}$} اور \lr{$\cos\frac{\pi}{4}$} کی قیمت \lr{$\frac{1}{\sqrt{2}}$} ہو تو درج ذیل مقدار معلوم کریں: \lr{$#2$}}}
\newcommand{\trigb}[2]{\uq{#1}{درج ذیل تکونی اکائی کو ثابت کریں:}\mline{#2}}
\newcommand{\ptsq}[2]{\uq{#1}{نقاط کے دو جوڑے دیے گئے ہیں۔ معلوم کریں کہ آیان نقاط سے بنے والے دو خطوط (الف) متوازی ہیں (ب) عمودی ہیں (ج) کچھ بھی نہیں ہیں \lr{$#2$}}}
\newcommand{\lineform}[2]{\uq{#1}{درج ذیل مساوات کو ڈھلوان قاطع شکل، دو قاطع شکل اور عام شکل میں تبدیل کریں۔}\mline{#2}}
\newcommand{\parperp}[2]{\uq{#1}{درج ذیل میں چیک کریں کہ آیا دو خطوط: (الف) متوازی ہیں (ب) عمودی ہیں (ج) نامتوازی ہیں اور نہ عمودی ہیں۔}\mline{#2}}
\newcommand{\dproof}[2]{\uq{#1}{درج ذیل الجبری جملے کے لیے سادہ استخراجی ثبوتوں کو وضع کریں اور ثابت کریں کہ بائیں طرف برابر ہے دائیں طرف کے: \lr{$#2$}}}
\newcommand{\jproof}[2]{\uq{#1}{ہر مرحلے کی وضاحت کرتے ہوئے درج ذیل کو ثابت کریں: \lr{$#2$}}}
\newcommand{\ptcoord}[2]{\uq{#1}{#2}}
\newcommand{\plotq}[2]{\uq{#1}{درج ذیل دو درجی اور سہ درجی تفاعل کا گراف بنائیں:}\mline{#2}}
\newcommand{\bisect}[2]{\uq{#1}{تصدیق کریں کہ درج ذیل پیمائشوں سے \lr{$\triangle ABC$} کے زاویوں کے ناصف ہم نقطہ ہوتے ہیں:}\mline{#2}}
\newcommand{\perpbis}[2]{\uq{#1}{دی گئی پیمائشوں سے \lr{$\triangle ABC$} بنائیں اور تصدیق کریں کہ مثلث کے عمودی ناصف ہم نقطہ ہوتے ہیں:}\mline{#2}}
\newcommand{\medq}[2]{\uq{#1}{درج ذیل پیمائشوں سے \lr{$\triangle LMN$} بنائیں اور تصدیق کریں کہ مثلث کے وسطانیے ہم نقطہ ہوتے ہیں:}\mline{#2}}

\begin{document}
\begin{center}
  {\Huge\bfseries The Hope Science Academy}\\[3pt]
  {\small Excellence is our standard not goal}\\[2pt]
  {\small 247-E-1, (Opposite LDA School) Johar Town, Lahore. PH: 0300-8194789, 0322-4157001}
\end{center}
\medskip\noindent
\begin{tabularx}{\linewidth}{|X|X|X|X|}
\hline
\textbf{Student Name:} & \textbf{Roll Num:} & \textbf{Class:} 9th & \textbf{Paper Code:} \\ \hline
\textbf{Subject Name:} Mathematics. & \textbf{Time Allowed:} & \textbf{Total Marks:} 3101 & \textbf{Exam Date:} 09-Oct-2026 \\ \hline
\multicolumn{4}{|l|}{\textbf{Exam Syllabus:} FULL BOOK} \\ \hline
\end{tabularx}
\bigskip
\medskip\noindent\textbf{1$\times$448=448}\hfill{\urdufont\beginR\textbf{سوال نمبر \lr{5}۔ مندرجہ ذیل سوالات کے تفصیلی جوابات تحریر کریں۔}\endR}\par
\uq{1}{اگر \lr{$a=\frac{2}{3},\ b=\frac{3}{2},\ c=\frac{5}{3}$} ہو تو خاصیت تقسیمی بلحاظ جمع کو ثابت کریں۔}
\uq{2}{درج ذیل کو مختصر کیجیے:}\mline{\frac{5.(25)^{n+1}-25.(5)^{2n}}{5.(5)^{2n+3}-(25)^{n+1}}}
\uq{3}{درج ذیل کو مختصر کیجیے:}\mline{\frac{3^{n}\times9^{n+1}}{3^{n-1}\times9^{n-1}}}
\uq{4}{درج ذیل کو مختصر کیجیے:}\mline{\frac{(16)^{x+1}+20(4^{2x})}{2^{x-3}\times8^{x+2}}}
\uq{5}{درج ذیل کو مختصر کیجیے:}\mline{\frac{5^{n+3}-6.5^{n+1}}{9\times5^{n}-2^{n}\times5^{n}}}
\uq{6}{اگر \lr{$x=3+\sqrt{8}$} ہو درج ذیل کی قیمت معلوم کریں:}\mline{x^{2}-\frac{1}{x^{2}}}
\uq{7}{اگر \lr{$\dfrac{8-3\sqrt{2}}{4+3\sqrt{2}}=p+q\sqrt{2}$} ہو تو ناطق اعداد \lr{p} اور \lr{q} معلوم کریں۔}
\uq{8}{اگر \lr{$x=3+\sqrt{8}$} ہو درج ذیل کی قیمت معلوم کریں:}\mline{x^{4}+\frac{1}{x^{4}}}
\uq{9}{درج ذیل کو مختصر کیجیے:}\mline{\frac{(25)^{\frac{3}{2}}\times(243)^{\frac{3}{5}}}{(16)^{\frac{5}{4}}\times(8)^{\frac{4}{3}}}}
\uq{10}{درج ذیل کو مختصر کیجیے:}\mline{\frac{54\times\sqrt[3]{(27)^{2x}}}{9^{x+1}+216\,(3^{2x-1})}}
\uq{11}{درج ذیل کو مختصر کیجیے:}\mline{\left(a^{\frac{1}{3}}+b^{\frac{2}{3}}\right)\times\left(a^{\frac{2}{3}}-a^{\frac{1}{3}}b^{\frac{2}{3}}+b^{\frac{4}{3}}\right)}
\uq{12}{درج ذیل کو مختصر کیجیے:}\mline{\sqrt{\frac{(216)^{\frac{2}{3}}\times(25)^{\frac{1}{2}}}{(0.04)^{\frac{-3}{2}}}}}
\uq{13}{شکل میں قائمۃ الزاویہ \lr{$\triangle ABC$} دکھائی گئی ہے جس میں \lr{$\overline{AC}$} کی لمبائی \lr{$(\sqrt{3}+\sqrt{5})$} سینٹی میٹر ہے۔ \lr{$\triangle ABC$} کا رقبہ \lr{$(1+\sqrt{15})\,cm^{2}$} ہے۔ ضلع \lr{$\overline{AB}$} کی لمبائی \lr{$(a\sqrt{3}+b\sqrt{5})\,cm$} کی صورت میں معلوم کریں، جہاں \lr{a} اور \lr{b} صحیح اعداد ہیں۔}
\tfig{\coordinate(A)at(0,0);\coordinate(B)at(0,2);\coordinate(C)at(4,0);
\draw(A)--(B)--(C)--cycle;\draw(0,.25)--(.25,.25)--(.25,0);
\node[left]at(A){$A$};\node[above]at(B){$B$};\node[right]at(C){$C$};
\node[below]at(2,0){$(\sqrt{3}+\sqrt{5})\,cm$};}
\uq{14}{ایک مستطیل کے اضلاع کی لمبایاں \lr{$2+\sqrt{18}\,m$} اور \lr{$\left(5-\frac{4}{\sqrt{2}}\right)$} میٹر ہیں۔ مستطیل کے رقبے کو \lr{$a+b\sqrt{2}$} کی شکل میں ظاہر کریں۔ جہاں \lr{a} اور \lr{b} صحیح اعداد ہیں۔}
\uq{15}{دو اعداد معلوم کریں جن کا مجموعہ \lr{68} ہے اور فرق \lr{22} ہے۔}
\uq{16}{2024 کے موسم گرما کے دوران لاہور کا موسم غیر معمولی طور پر گرم تھا۔ ٹی وی نیوز نے درجہ حرارت 48 ڈگری سینٹی گریڈ تک بتایا۔ فارمولہ استعمال کرتے ہوئے، \lr{$\left({}^{o}F=\frac{9}{5}{}^{o}C+32\right)$} درجہ حرارت کو فارن ہائیٹ پیمانے کے طور پر معلوم کریں۔}
\uq{17}{باپ بیٹے کی عمروں کا مجموعہ \lr{72} سال ہے۔ چھ سال پہلے باپ کی عمر بیٹے کی عمر سے دو گنا تھی۔ چھ سال پہلے بیٹے کی عمر کیا تھی؟}
\uq{18}{میرا ایک کھلونا \lr{1520} روپے میں بیچتی ہے۔ \lr{15\%} منافع حاصل کرنے کے لیے قیمت فروخت کیا ہوگی؟}
\uq{19}{طیب کی سالانہ آمدنی \lr{9,60,000} روپے ہے۔ جبکہ مستثنیٰ (چھوٹ) رقم \lr{1,30,000} روپے ہے۔ اسے \lr{0.75\%} کی شرح سے کتنا ٹیکس ادا کرنا پڑے گا؟}
\uq{20}{\lr{3,75,000} روپے پر کمپاؤنڈ مارک اپ معلوم کریں۔ ایک سال کے لیے \lr{14\%} سالانہ مرکب کی شرح سے۔}
\uq{21}{تین یکے بعد دیگرے طاق عدد کا مجموعہ \lr{51} ہے۔ تین اعداد معلوم کریں۔}
\uq{22}{عبداللہ نے \lr{96} گیندیں اٹھائیں اور انہیں دو بالٹیوں میں ڈالا۔ ایک بالٹی میں دوسری بالٹی سے اٹھائیس زیادہ گیندیں ہیں۔ ہر بالٹی میں کتنی گیندیں تھیں؟}
\uq{23}{سلمیٰ نے \lr{3,50,000} روپے کی سرمایہ کاری کی۔ ایک بینک میں، جس نے \lr{$7\frac{1}{4}\%$} سالانہ کی شرح سے سادہ منافع ادا کیا۔ 2 سال کے بعد، شرح کو بڑھا کر \lr{8\%} سالانہ کر دیا گیا۔ 7 سال کے آخر میں اس کے پاس موجود رقم معلوم کریں۔}
\uq{24}{حمائل نے \lr{6590} روپے میں ایک سائیکل خریدی اور اسے \lr{6850} روپے میں فروخت کیا۔ منافع کا فیصد معلوم کریں۔}
\uq{25}{عمیر نے ایک کتاب \lr{850} روپے میں خریدی۔ \lr{720} روپے میں فروخت کیا۔ اس کے نقصان کا فیصد کیا تھا؟}
\uq{26}{سلیم، ندیم، اور تنویر نے ایک کاروبار سے \lr{4,50,000} روپے کا منافع کمایا۔ اس کاروبار میں ان کی سرمایہ کاری کا تناسب \lr{4:7:14} ہے، تو ہر شخص کا منافع معلوم کریں۔}
\uq{27}{اگر \lr{6400} روپے پر سادہ منافع \lr{12} سال کے لیے \lr{3840} روپے ہے۔ منافع کی شرح معلوم کریں۔}
\uq{28}{درج ذیل میں \lr{x} کی قیمت معلوم کریں: \lr{$\log_{4}x=\frac{3}{2}$}}
\uq{29}{درج ذیل میں \lr{x} کی قیمت معلوم کریں: \lr{$\log_{2}1024=x$}}
\uq{30}{\lr{x} کی قیمت معلوم کیجیے: \lr{-3.434}}
\uq{31}{\lr{x} کی قیمت معلوم کیجیے: \lr{-1.5726}}
\uq{32}{\lr{x} کی قیمت معلوم کیجیے: \lr{4.3561}}
\uq{33}{\lr{x} کی قیمت معلوم کیجیے: \lr{log x = 2.0184}}
\uq{34}{\lr{x} کی قیمت معلوم کیجیے: \lr{-2.1234}}
\uq{35}{کیلکولیٹر کا استعمال کیے بغیر درج ذیل کی قیمت معلوم کیجیے: \lr{$\frac{1}{3}\log_{3}8-\log_{3}18$}}
\uq{36}{کیلکولیٹر کا استعمال کیے بغیر درج ذیل کی قیمت معلوم کیجیے: \lr{$\frac{1}{3}\log_{4}64+2\log_{5}25$}}
\uq{37}{کیلکولیٹر کا استعمال کیے بغیر درج ذیل کی قیمت معلوم کیجیے: \lr{$\log_{3}12+\log_{3}0.25$}}
\uq{38}{درج ذیل مساوات میں \lr{x} کی قیمت معلوم کیجیے: \lr{$(81)^{x}=(243)^{x+2}$}}
\uq{39}{درج ذیل مساوات میں \lr{x} کی قیمت معلوم کیجیے: \lr{$\left(\frac{1}{27}\right)^{x-6}=27$}}
\uq{40}{درج ذیل مساوات میں \lr{x} کی قیمت معلوم کیجیے: \lr{$\log_{2}(x+1)-\log_{2}(x-4)=2$}}
\uq{41}{لوگارتھم جدول کی مدد سے درج ذیل کی قیمت معلوم کیجیے: \lr{$\dfrac{3.68\times4.21}{5.234}$}}
\uq{42}{لوگارتھم جدول کی مدد سے درج ذیل کی قیمت معلوم کیجیے: \lr{$4.67\times2.11\times2.397$}}
\uq{43}{لوگارتھم جدول کی مدد سے درج ذیل کی قیمت معلوم کیجیے: \lr{$\dfrac{(20.46)^{2}\times(2.4122)}{754.3}$}}
\uq{44}{لوگارتھم جدول کی مدد سے درج ذیل کی قیمت معلوم کیجیے: \lr{$\dfrac{\sqrt[3]{9.364}\times21.64}{3.21}$}}
\uq{45}{زلزلوں کی شدت کی پیمائش کا کلیہ \lr{$M=\log_{10}\left(\frac{A}{A_{o}}\right)$} ہے۔ اگر \lr{10,000} طول موج \lr{(A)} اور \lr{10} طول موج کا حوالہ \lr{(A$_{o}$)} ہو تو زلزلے کی شدت کتنی ہوگی؟}
\uq{46}{عبداللہ نے ایک بچت اسکیم میں \lr{100,000} روپے کی سرمایہ کاری کی اور سالانہ \lr{5\%} کی شرح سے منافع حاصل کیا تاکہ اس سرمایہ کاری کی کل مالیت \lr{t} سالوں کے بعد \lr{y} روپے ہو جائے۔ اے ایک مساوات \lr{$y=100,000\,(1.05)^{t},\ t\geq0$} کی شکل میں ظاہر کیا گیا ہے۔ معلوم کریں کہ کتنے سالوں بعد سرمایہ کاری دوگنا ہو جائے گی۔}
\uq{47}{حوریہ ایک ایسے پہاڑ پر چڑھ رہی ہے جہاں درجہ حرارت ہر \lr{100} میٹر اونچائی بڑھنے پر \lr{3\%} (یا \lr{0.97} کا جزو) کم ہو جاتا ہے۔ ابتدائی درجہ حرارت \lr{(T$_{i}$)} سطح سمندر پر \lr{20$^{o}$C} ہے۔ کلیہ \lr{$T=T_{i}\times0.97^{\frac{h}{100}}$} کو استعمال کرتے ہوئے \lr{500} میٹر کی اونچائی پر درجہ حرارت معلوم کریں۔}
\uq{48}{ڈیسیبل اسکیل کا کلیہ \lr{$L=40\log_{10}\left(\frac{I}{I_{o}}\right)$} آواز کی شدت کی پیمائش کرتا ہے۔ اگر کسی آواز کی شدت \lr{(I)} حوالہ شدت \lr{(I$_{o}$)} کا \lr{$10^{6}$} گنا ہو تو ڈیسیبل میں آواز کی سطح کیا ہوگی؟}
\uq{49}{درج ذیل میں \lr{x} کی قیمت معلوم کریں: \lr{$\left(\frac{1}{32}\right)^{2x}=64$}}
\uq{50}{درج ذیل کی قیمت لوگارتھم جدول کی مدد سے معلوم کیجیے: \lr{319.8$\times$3.543}}
\uq{51}{رج ذیل کی قیمت لوگارتھم جدول کی مدد سے معلوم کیجیے: \lr{$\dfrac{36.12\times750.9}{113.2\times9.98}$}}
\uq{52}{سال \lr{2016} میں ایک شہر کی آبادی \lr{22} ملین تھی اور ہر سال \lr{2.5\%} کی شرح سے بڑھ رہی تھی۔ تفاعل \lr{$p(t)=22(1.025)^{t}$} ملین میں \lr{2016} کے بعد \lr{t} سالوں میں آبادی معلوم کرتا ہے۔ دیے گئے تفاعل کا استعمال کرتے ہوئے معلوم کریں کہ کتنے سالوں میں آبادی \lr{35} ملین تک پہنچ جائے گی۔}
\uq{53}{یونیورسل سیٹ پر غور کریں \lr{U = \{x$\mid$ x ضعف ہے 2 کا $\wedge\ 0<x\leq30$\}}، \lr{A = \{x$\mid$ x ضعف ہے 6 کا\}} اور \lr{B = \{x$\mid$ x ضعف ہے 8 کا\}} ہو تو وین اشکال بنائیں}
\uq{54}{فرض کریں \lr{U = \{x$\mid$ x} ایک صحیح عدد ہے \lr{$\wedge\ 0<x\leq150$\}} \lr{G = \{x$\mid$ x $=2^{m}$} کے لیے \lr{m} ایک صحیح عدد\lr{\}} اور \lr{H = \{x$\mid$ x} ایک مربع ہے\lr{\}} ہو تو سیٹ \lr{G} اور \lr{H} کے تمام ارکان کو اندراجی طریقہ میں لکھیں۔}
\uq{55}{فرض کریں \lr{U = \{x$\mid$ x} ایک صحیح عدد ہے \lr{$\wedge\ 0<x\leq150$\}} \lr{G = \{x$\mid$ x $=2^{m}$} کے لیے \lr{m} ایک صحیح عدد\lr{\}} اور \lr{H = \{x$\mid$ x} ایک مربع ہے\lr{\}} ہو تو \lr{$G\cup H$} معلوم کریں۔}
\uq{56}{فرض کریں \lr{U = \{x$\mid$ x} ایک صحیح عدد ہے \lr{$\wedge\ 0<x\leq150$\}} \lr{G = \{x$\mid$ x $=2^{m}$} کے لیے \lr{m} ایک صحیح عدد\lr{\}} اور \lr{H = \{x$\mid$ x} ایک مربع ہے\lr{\}} ہو تو \lr{$G\cap H$} معلوم کریں۔}
\uq{57}{فرض کریں \lr{P = \{x$\mid$ x} ایک مفرد عدد ہے \lr{$\wedge\ 0<x\leq20$\}} اور \lr{Q = \{x$\mid$ x $\leq20\ \wedge$ x 210} کا ایک مقسوم علیہ ہے\lr{\}} ہو تو \lr{$P\cap Q$} معلوم کریں۔}
\uq{58}{فرض کریں \lr{P = \{x$\mid$ x} ایک مفرد عدد ہے \lr{$\wedge\ 0<x\leq20$\}} اور \lr{Q = \{x$\mid$ x $\leq20\ \wedge$ x 210} کا ایک مقسوم علیہ ہے\lr{\}} ہو تو \lr{$P\cup Q$} معلوم کریں۔}
\uq{59}{اگر \lr{U=\{1,2,3,...,20\}} اور \lr{A=\{1,3,5,...,19\}}، ہو تو درج ذیل کو ثابت کریں: \lr{$A\cup A'=U$}}
\uq{60}{اگر \lr{U=\{1,2,3,...,20\}} اور \lr{A=\{1,3,5,...,19\}}، ہو تو درج ذیل کو ثابت کریں: \lr{$A\cap A'=\phi$}}
\uq{61}{\lr{55} طلبہ کی ایک کلاس میں \lr{34} کرکٹ کھیلنا پسند کرتے ہیں۔ اور \lr{30} ہاکی کھیلنا پسند کرتے ہیں۔ نیز ہر طالب علم کم از کم دو میں سے ایک کھیل کھیلنا پسند کرتا ہے۔ کتنے طلبہ دونوں کھیل کھیلنا پسند کرتے ہیں؟}
\uq{62}{\lr{500} ملازمین کے گروہ میں \lr{250} اردو بول سکتے ہیں \lr{150} انگریزی بول سکتے ہیں۔ \lr{50} پنجابی بول سکتے ہیں۔ \lr{40} اردو اور انگریزی بول سکتے ہیں۔ \lr{30} انگریزی اور پنجابی دونوں بول سکتے ہیں۔ اور \lr{10} اردو اور پنجابی بول سکتے ہیں۔ کتنے لوگ تینوں زبانیں بول سکتے ہیں؟}
\uq{63}{کھیلوں کے مقابلوں میں \lr{19} لوگ نیلی شرٹس پہنتے ہیں۔ \lr{15} سبز قمیض پہنتے ہیں۔ \lr{3} نیلی اور سبز قمیضیں پہنتے ہیں۔ \lr{4} ٹوپی اور نیلی قمیض پہنتے ہیں۔ اور \lr{2} ٹوپی اور سبز قمیض پہنتے ہیں۔ نیلی یا سبز قمیض یا ٹوپی والے لوگوں کی کل تعداد \lr{34} ہے۔ کتنے لوگ ٹوپی پہنے ہوئے ہیں؟}
\uq{64}{ایک تربیتی نشست میں \lr{17} شرکا کے پاس لیپ ٹاپ، \lr{11} کے پاس ٹیبلیٹ، \lr{9} کے پاس لیپ ٹاپ اور ٹیبلیٹ، \lr{6} کے پاس لیپ ٹاپ اور کتب ہیں۔ \lr{8} شرکاء کے پاس تینوں اشیاء ہیں۔ لیپ ٹاپ، ٹیبلیٹ یا کتب کے ساتھ شرکا کی کل تعداد \lr{35} ہے۔ کتنے شرکا کے پاس کتب ہیں؟}
\uq{65}{\mallemp\ \lr{$(A'\cup B')\cap C$} معلوم کریں۔}
\uq{66}{\mallemp\ \lr{$n\{A\cap(B^{c}\cap C^{c})\}$} معلوم کریں۔}
\uq{67}{\sportsch\ کتنے طلبہ تینوں کھیل کھیلتے ہیں؟}
\uq{68}{\sportsch\ تمام کھیلوں میں شرکت کی تقسیم کو ظاہر کرنے کے لیے وین اشکال بنائیں۔}
\uq{69}{\foodsurvey\ کم از کم کتنے لوگوں کو نہاری، بریانی یا قورمہ پسند ہے؟}
\uq{70}{\foodsurvey\ کتنے لوگوں کو نہاری، بریانی یا قورمہ پسند نہیں آیا؟}
\uq{71}{\foodsurvey\ کتنے لوگ درجہ بالا کھانوں میں سے صرف ایک کو پسند کرتے ہیں، نہاری، بریانی یا قورمہ؟}
\uq{72}{\foodsurvey\ وین اشکال بنائیں۔}
\uq{73}{ملحقہ وین ڈایاگرام پر غور کریں جو دو غیر خالی سیٹوں \lr{A} اور \lr{B} کو ظاہر کرتی ہے۔ سیٹ \lr{A} اور \lr{B} کا یونین معلوم کریں۔}
\pfig{fig73}{5.5cm}
\uq{74}{ملحقہ وین ڈایاگرام پر غور کریں جو دو غیر خالی سیٹوں \lr{A} اور \lr{B} کو ظاہر کرتی ہے۔ صرف سیٹ \lr{A} کے وہ تمام ارکان معلوم کریں جو سیٹ \lr{A} میں موجود نہ ہوں۔}
\pfig{fig74}{5.5cm}
\uq{75}{ملحقہ وین ڈایاگرام پر غور کریں جو دو غیر خالی سیٹوں \lr{A} اور \lr{B} کو ظاہر کرتی ہے۔ سیٹ \lr{A} اور \lr{B} میں مشترک ارکان کی تعداد معلوم کریں۔}
\pfig{fig75}{5.5cm}
\uq{76}{\itvenn\ کتنے طلبہ نے اپلائیڈ روبوٹکس کورس میں داخلہ لیا؟}
\pfig{fig76}{5cm}
\uq{77}{ملحقہ وین شکل پر غور کریں جس میں ایک آئی ٹی ادارے میں مختلف کورسز میں داخلہ لینے والے طلبہ کو ظاہر کیا گیا ہے۔ \lr{U=\{}آئی ٹی ادارے میں داخلہ لینے والے طلبہ\lr{\}}، \lr{A= \{}اپلائیڈ روبوٹکس میں داخلہ لینے والے طلبہ\lr{\}}، \lr{G= \{}گیم ڈویلپمنٹ میں داخلہ لینے والے طلبہ\lr{\}} \lr{W=\{} ویب ڈیزائننگ میں داخلہ لینے والے طلبہ۔ گیم ڈویلپمنٹ میں داخلہ لینے والے طلبہ\lr{\}} کی کل تعداد معلوم کریں۔}
\pfig{fig76}{5cm}
\uq{78}{\itvenn\ گیم ڈویلپمنٹ اور ویب ڈیزائننگ کورس میں کتنے طلبہ نے داخلہ لیا؟}
\pfig{fig76}{5cm}
\uq{79}{\itvenn\ ان تمام طلبہ کی تعداد معلوم کریں جنہوں نے ویب ڈویلپمنٹ کے کورس میں داخلہ لیا لیکن اپلائیڈ روبوٹکس کے کورس میں داخلہ نہیں لیا۔}
\pfig{fig76}{5cm}
\uq{80}{\itvenn\ آئی ٹی کے ادارے میں کتنے طلبہ داخل ہیں؟}
\pfig{fig76}{5cm}
\uq{81}{\itvenn\ تینوں کورسز میں کتنے طلبہ نے داخلہ لیا؟}
\pfig{fig76}{5cm}
\uq{82}{ایک اسپورٹس کلب میں \lr{98} سیکنڈری سکول کے طلبہ ہیں۔ \lr{58} طلبہ سوئمنگ کلب میں شامل ہوئے اور \lr{50} نے رسہ کشی کلب میں شمولیت اختیار کی۔ دونوں کھیلوں میں کتنے طلبہ نے حصہ لیا؟}
\uq{83}{سلیم جو ایک سکول ٹیچر ہے۔ اس کے گھر میں ایک چھوٹی لائبریری ہے۔ جس میں \lr{150} کتب ہیں۔ ان کے پاس ان کتب کی دو اہم اقسام ہیں۔ اسلامی اور سائنسی۔ انہوں نے \lr{70} کتب کو اسلام کتب اور \lr{90} کتب کو سائنس کی کتب میں درجہ بندی کیا۔ \lr{15} کتب ایسی ہیں جو نہ تو اسلامی اور نہ ہی سائنسی کتب کے زمرے میں آتی ہیں۔ کتنی کتب اسلامی اور سائنسی دونوں زمروں میں درجہ بندی کی گئی ہیں؟}
\uq{84}{ایک کالج میں \lr{45} اساتذہ ریاضی یا فزکس یا کیمسٹری پڑھاتے ہیں۔ یہاں مختلف مضامین پڑھانے والے اساتذہ کے بارے معلومات دی گئی ہے: \lr{18} ریاضی پڑھاتے ہیں، \lr{12} فزکس پڑھاتے ہیں، \lr{8} کیمسٹری پڑھاتے ہیں۔ \lr{6} ریاضی اور فزکس دونوں پڑھاتے ہیں۔ \lr{4} فزکس اور کیمسٹری دونوں پڑھاتے ہیں۔ \lr{2} ریاضی اور کیمسٹری دونوں پڑھاتے ہیں۔ کتنے اساتذہ تینوں مضامین پڑھاتے ہیں؟}
\uq{85}{\shopsurvey\ کتنے صارفین نے کم از کم اور مصنوعات خریدی: گارمنٹس، سنگھار یا الیکٹرونکس؟}
\uq{86}{\shopsurvey\ کتنے صارفین نے صرف ایک مصنوعات خریدی: گارمنٹس، سنگھار یا الیکٹرونکس؟}
\uq{87}{\shopsurvey\ کتنے صارفین نے تین میں سے کوئی بھی مصنوعات نہیں خریدی؟}
\uq{88}{فرض کریں \lr{$f(x)=ax+b+1$} جب کہ \lr{a} اور \lr{b} مستقل عدد ہیں۔ اگر \lr{$f(3)=8$} اور \lr{$f(6)=14$} ہو تو \lr{a} اور \lr{b} کی قیمتیں معلوم کریں۔}
\uq{89}{فرض کریں۔ \lr{$g(x)=ax+b+5$} جب کہ \lr{a} اور \lr{b} مستقل اعداد ہیں۔ اگر \lr{$g(-1)=0$} اور \lr{$g(2)=10$} ہو تو \lr{a} اور \lr{b} کی قیمتیں معلوم کریں۔}
\uq{90}{فرض کریں \lr{$f(x)=cx^{2}+d$} جب کہ \lr{c} اور \lr{d} مستقل اعداد ہیں۔ اگر \lr{$f(1)=6$} اور \lr{$f(-2)=10$} ہو تو \lr{c} اور \lr{d} کی قیمتیں معلوم کریں۔}
\uq{91}{فرض کریں \lr{$c_{1},c_{2},c_{3}$} تین بچے ہوں اور \lr{$m_{1},m_{2}$} دو آدمی ایسے ہوں کے دونوں \lr{$c_{1},c_{2}$} کا باپ \lr{m$_{1}$} ہو اور \lr{$c_{3}$} کا باپ \lr{m$_{2}$} ہو۔ ربط معلوم کریں۔ \lr{\{(بیٹا,باپ)\}}}
\uq{92}{فرض کریں: \lr{$f(x)=ax+b+3$} جب کہ \lr{a} اور \lr{b} مستقل اعداد ہیں۔ اگر \lr{$f(1)=4$} اور \lr{$f(5)=9$} ہو تو \lr{a} اور \lr{b} کی قیمتیں معلوم کریں۔}
\uq{93}{وین اشکال کا استعمال کرتے ہوئے درج ذیل کو ثابت کریں: \lr{$A-B=A\cup B'$}}
\uq{94}{وین اشکال کا استعمال کرتے ہوئے درج ذیل کو ثابت کریں: \lr{$(A-B)'\cap B=B$}}
\uq{95}{\abprop\ \lr{$A=\{1,2,3,4\},\ B=\{3,4,5,6,7,8\},\ C=\{5,6,7,9,10\}$}}
\uq{96}{\abprop\ \lr{$A=\phi,\ B=\{0\},\ C=\{0,1,2\}$}}
\uq{97}{\abprop\ \lr{$A=N,\ B=Z,\ C=Q$}}
\uq{98}{درج ذیل سیٹوں کے لیے ڈی مورگن کے قوانین کو ثابت کریں: \lr{$U=\{1,2,3,....,20\},\ A=\{2,4,6,....,20\}$} اور \lr{$B=\{1,3,5,....,19\}$}}
\uq{99}{یونین اور تقاطع کی مناسب خصوصیات کی مدد سے، درج ذیل نتائج اخذ کریں۔ \lr{$A\cap(A\cup B)=A\cup(A\cap B)$}}
\uq{100}{یونین اور تقاطع کی مناسب خصوصیات کی مدد سے، درج ذیل نتائج اخذ کریں۔ \lr{$A\cup(A\cap B)=A\cap(A\cup B)$}}
\uq{101}{فرض کریں \lr{$f(x)=ax+b$} جبکہ \lr{a} اور \lr{b} مستقل اعداد ہیں۔ اگر \lr{$f(-2)=3$} اور \lr{$f(4)=10$} ہو تو \lr{a} اور \lr{b} کی قیمتیں معلوم کریں۔}
\uq{102}{فرض کریں \lr{$g(x)=mx^{2}+n$} جب کہ \lr{m} اور \lr{n} مستقل اعداد ہیں۔ اگر \lr{g(4)=20} اور \lr{g(0)=5} ہو تو \lr{m} اور \lr{n} کی قیمتیں معلوم کریں۔}
\uq{103}{\examtxt\ کتنے طلبہ نے دونوں ٹیسٹوں میں سے کوئی بھی پاس نہیں کیا؟}
\uq{104}{\examtxt\ کتنے طلبہ نے سائنس کا امتحان پاس کیا لیکن ریاضی کا امتحان نہیں؟}
\uq{105}{\examtxt\ سائنس کے امتحان میں کتنے طلبہ فیل ہوئے؟}
\uq{106}{فرض کریں۔ \lr{$U=\{1,2,3,4,5,6,7,8,9,10\},\ A=\{2,4,6,8,10\},\ B=\{1,2,3,4,5\}$} اور \lr{$C=\{1,3,5,7,9\}$} ہو تو درج ذیل سیٹ کے ارکان کی فہرست بنائیں۔ \lr{$A'\cup C'$}}
\uq{107}{فرض کریں۔ \lr{$U=\{1,2,3,4,5,6,7,8,9,10\},\ A=\{2,4,6,8,10\},\ B=\{1,2,3,4,5\}$} اور \lr{$C=\{1,3,5,7,9\}$} ہو تو درج ذیل سیٹ کے ارکان کی فہرست بنائیں۔ \lr{$A'\cup C$}}
\uq{108}{ایک شاپنگ مال میں \lr{1} سے \lr{100} لیبل والی مختلف اقسام کے \lr{100} اشیا ہیں۔ جو یونیورسل سیٹ \lr{U} کو ظاہر کرتے ہیں۔ اشیا کی درجہ بندی درج ذیل ہے: سیٹ \lr{A:} الیکٹرونکس \lr{1} سے \lr{30} تک کے لیبل والی \lr{30} اشیا پر مشتمل ہے۔ سیٹ \lr{B:} کپڑوں میں \lr{31} سے \lr{55} تک کے لیبل والی \lr{25} اشیا شامل ہیں۔ سیٹ \lr{C:} سنگھار کی اشیا، جس میں \lr{76} سے \lr{100} تک کے لیبل والی \lr{25} اشیا شامل ہیں۔ ہر سیٹ کو اندراجی طریقہ میں لکھیں اور تینوں سیٹوں کا یونین معلوم کریں۔}
\uq{109}{\examtxt\ کتنے طلبہ نے ریاضی یا سائنس کا ٹیسٹ پاس کیا؟}
\fzA{110}{2x^{2}+11x+15}
\fzA{111}{3y^{2}-11y+6}
\fzA{112}{2y^{2}-5y+2}
\fzA{113}{4z^{2}-11z+6}
\fzA{114}{6+7x-3x^{2}}
\fzA{115}{2x^{2}+17x+26}
\fzB{116}{4x^{4}+81y^{4}}
\fzB{117}{a^{4}+64b^{4}}
\fzB{118}{x^{4}+4x^{2}+16}
\fzB{119}{x^{4}-14x^{2}+1}
\fzB{120}{x^{4}-30x^{2}y^{2}+9y^{4}}
\fzB{121}{x^{4}+11x^{2}y^{2}+y^{4}}
\fzB{122}{(x+4)(x+3)(x+2)(x+1)+1}
\fzB{123}{(x-1)(x-4)(x-7)(x+2)+17}
\fzB{124}{(2x^{2}+7x+3)(2x^{2}+7x+5)+1}
\fzB{125}{(3x^{2}+5x+3)(3x^{2}+5x+5)-3}
\fzB{126}{(x+1)(x+2)(x+3)(x+6)-3x^{2}}
\fzB{127}{(x+1)(x-1)(x+2)(x-2)+13x^{2}}
\fzA{128}{a^{4}+64}
\fzA{129}{(x^{2}+5x+4)(x^{2}+5x+6)-3}
\fzA{130}{(x-2)(x+2)(x+1)(x-4)+2x^{2}}
\fzA{131}{8x^{3}+60x^{2}+150x+125}
\fzA{132}{x^{4}+x^{2}+25}
\uq{133}{تجزی کریں \lr{$x^{4}+y^{4}$}}
\hcffac{134}{a^{3}+2a^{2}-3a,\ 2a^{3}+5a^{2}-3a}
\hcffac{135}{t^{2}+3t-4,\ t^{2}+5t+4,\ t^{2}-1}
\hcffac{136}{x^{2}+15x+56,\ x^{2}+5x-24,\ x^{2}+8x}
\uq{137}{درج ذیل جملے کا عاد اعظم بذریعہ تقسیم معلوم کریں: \lr{$27x^{3}+9x^{2}-3x-9,\ 3x-2$}}
\hcfdiv{138}{x^{3}-9x^{2}+21x-15,\ x^{2}-4x+3}
\hcfdiv{139}{2x^{3}+2x^{2}+2x+2,\ 6x^{3}+12x^{2}+6x+12}
\uq{140}{درج ذیل جملے کا عاد اعظم بذریعہ تقسیم معلوم کریں: \lr{$2x^{3}-4x^{2}+6x,\ x^{3}-2x,\ 3x^{2}-6x$}}
\uq{141}{درج ذیل جملے کا ذواضعاف اقل بذریعہ تجزی معلوم کریں: \lr{$x^{4}-16,\ x^{3}-4x$}}
\uq{142}{درج ذیل جملے کا ذواضعاف اقل بذریعہ تجزی معلوم کریں:}\mline{16-4x^{2},\ x^{2}+x-6,\ 4-x^{2}}
\uq{143}{دو کثیر رقمی جملوں کا عاد اعظم \lr{$y-7$} اور ذواضعاف اقل \lr{$y^{3}-10y^{2}+11y+70$} ہے۔ اگر ایک کثیر رقمی جملہ \lr{$y^{2}-5y-14$} ہو تو دوسرا کثیر رقمی جملہ معلوم کریں۔}
\uq{144}{دو کثیر رقمی جملوں \lr{p(x)} اور \lr{q(x)} کا ذواضعاف اقل اور عاد اعظم بالترتیب \lr{$36x^{3}(x+a)(x^{3}-a^{3})$} اور \lr{$x^{2}(x-a)$} ہے۔ اگر \lr{$p(x)=4x^{2}(x^{2}-a^{2})$} ہو تو \lr{q(x)} اور \lr{q(x)} معلوم کریں۔}
\uq{145}{دو کثیر رقمی جملوں کا عاد اعظم اور ذواضعاف اقل بالترتیب \lr{$(x+a)$} اور \lr{$12x^{2}(x+a)(x^{2}-a^{2})$} ہو تو دونوں جملوں کا حاصل ضرب معلوم کریں۔}
\uq{146}{کثیر رقمی جملوں \lr{$x^{2}-27,\ x^{2}+6x-27,\ x^{2}-9$} کا عاد اعظم بذریعہ تجزی معلوم کریں۔}
\uq{147}{کثیر رقمی جملوں \lr{$6x^{3}-17x^{2}-5x+6$} اور \lr{$6x^{3}-5x^{2}-3x+2$} کا عاد اعظم بذریعہ تقسیم معلوم کریں۔}
\uq{148}{کثیر رقمی جملوں \lr{$x^{2}-3x+2,\ x^{2}-1$} اور \lr{$x^{2}-5x+4$} کا ذواضعاف اقل معلوم کریں۔}
\uq{149}{دو کثیر رقمی جملوں کا ذواضعاف اقل اور عاد اعظم بالترتیب \lr{$x^{3}-10x^{2}+11x+70$} اور \lr{$x-7$} ہے۔ اگر ایک کثیر رقمی جملہ \lr{$x^{2}-12x+35$} ہو تو دوسرا کثیر رقمی جملہ معلوم کریں۔}
\uq{150}{دو کثیر رقمی جملوں \lr{$x^{2}y+xy^{2}$} اور \lr{$x^{2}+xy$} کا ذواضعاف اقل \lr{$xy(x+y)$} ہے۔ ان کا عاد اعظم معلوم کریں۔}
\uq{151}{درج ذیل جملے کا جذر المربع بذریعہ تجزی معلوم کریں۔ \lr{$200t^{2}-120t+18$}}
\uq{152}{درج ذیل جملے کا جذر المربع بذریعہ تجزی معلوم کریں۔ \lr{$40x^{2}+120x+90$}}
\uq{153}{درج ذیل جملے کا جذر المربع بذریعہ تقسیم معلوم کریں: \lr{$4x^{4}-28x^{3}+37x^{2}+42x+9$}}
\uq{154}{برقی میدان میں پوٹینشل توانائی \lr{V(x)} فاصلہ \lr{x} کے سہ درجی تفاعل کے طور پر مختلف ہوتی ہے \lr{$V(x)=2x^{3}-6x^{2}+4x$} معلوم کریں کہ پوٹینشل توانائی کہاں صفر ہے۔}
\uq{155}{ساختی انجینئرنگ میں شہتیر کا انحراف \lr{Y(x)} دیا گیا ہے: \lr{$Y(x)=2x^{2}-8x+6$} یہ مساوات شہتیر کے ساتھ کسی بھی نقطہ \lr{x} پر عمودی انحراف دیتی ہے۔ صفر انحراف کے نقاط معلوم کریں۔}
\sqrdiv{156}{121x^{4}-198x^{3}-183x^{2}+216x+144}
\sqrdiv{157}{x^{4}-10x^{3}y+27x^{2}y^{2}-10xy^{3}+y^{4}}
\sqrdiv{158}{x^{4}-10x^{3}y+27x^{2}y^{2}-10xy^{3}+y^{4}}
\uq{159}{\lr{x} ہزار روپے کی سرمایہ کاری کے بعد ایک سرمایہ کار کی واپسی \lr{R(x)} روپے میں دو درجی جملہ دیا گیا ہے: \lr{$R(x)=-x^{2}+6x-8$} تجزی کریں اور سرمایہ کاری کی سطحیں معلوم کریں جس کے نتیجے میں صفر کی واپسی ہوتی ہے۔}
\uq{160}{کسی شے کی \lr{x} اکائیاں فروخت کرنے سے روپوں میں کمپنی کا منافع \lr{P(x)} سہ درجی جملے کے ذریعے دیا گیا ہے: \lr{$P(x)=x^{3}-15x^{2}+75x-125$} مساوی نقطہ معلوم کریں جہاں منافع صفر ہو۔}
\uq{161}{\lr{$36x^{4}-36x^{2}+9$} کا جذر المربع معلوم کریں۔}
\uq{162}{\lr{$x^{4}-12x^{3}+42x^{2}-36x+9$} کا جذر المربع معلوم کریں۔}
\uq{163}{کسی پرزے کی تیاری کے لیے لاگت کا تفاعل اس طرح ماڈل دیا گیا ہے: \lr{$C(x)=5x^{2}-25x+30$} یہاں \lr{x} جزو کی چوڑائی اور \lr{C(x)} قیمت ہے۔ \lr{x} کی قیمت معلوم کریں یہاں \lr{C(x)} کم سے کم ہو۔}
\uq{164}{کمپنی کے منافع \lr{P(x)} کو دو درجی مساوات سے ظاہر کیا گیا ہے۔ \lr{$P(x)=-5x^{2}+50x-120$} یہاں \lr{x} تیار شدہ اکائیوں کی تعداد کو ظاہر کرتا ہے اور \lr{P(x)} منافع کو ڈالر میں ظاہر کرتا ہے۔ معلوم کریں کہ منافع کو زیادہ کرنے کے لیے کتنے یونٹ تیار کرنے چاہییں۔}
\uq{165}{کیوبک پوٹینشل میں حرکت کرنے والے جسم کی پوٹینشل توانائی \lr{U(x)} کو اس طرح ظاہر کیا جاتا ہے: \lr{$U(x)=x^{3}-6x^{2}+12x-8$} جملہ کی تجزی کرتے ہوئے ان نقاط کو معلوم کریں جہاں توانائی کم سے کم ہو جائے۔}
\uq{166}{درج ذیل جملے کی تجزی کریں: \lr{$(x+4)(x+3)(x+5)(x+6)-360$}}
\uq{167}{درج ذیل جملے کی تجزی کریں: \lr{$36+(x^{2}+6x+3)(x^{2}+6x-9)$}}
\uq{168}{ذواضعاف اقل اور عاد اعظم معلوم کریں: \lr{$4x^{3}+12x^{2},\ 8x^{2}+16x$}}
\uq{169}{ذواضعاف اقل اور عاد اعظم معلوم کریں: \lr{$x^{3}+3x^{2}-4x,\ x^{2}-x-6$}}
\uq{170}{حوریہ اپنے قرض کی کل لاگت کا تجزیہ کر رہی ہے، جس کا ماڈل \lr{$C(x)=x^{2}-8x+15$} ہے جہاں \lr{x} سالوں کی تعداد کو ظاہر کرتا ہے۔ حوریہ کے قرض کی ادائیگی کی بہترین مدت کیا ہے؟}
\uq{171}{درج ذیل مساوات کو حل کریں اور عددی خط پر ظاہر کریں: \lr{$\dfrac{2x-1}{3}-\dfrac{3x}{4}=\dfrac{5}{6}$}}
\uq{172}{درج ذیل مساوات کو حل کریں اور عددی خط پر ظاہر کریں: \lr{$\dfrac{-5x}{10}=9-\dfrac{10}{5}x$}}
\uq{173}{درج ذیل غیر مساوات کو حل کریں اور عددی خط پر ظاہر کریں: \lr{$\dfrac{5}{3}x-\dfrac{3}{4}<\dfrac{-1}{12}$}}
\uq{174}{درج ذیل غیر مساوات کو حل کریں اور عددی خط پر ظاہر کریں: \lr{$\dfrac{1}{4}x-\dfrac{1}{2}\leq-1+\dfrac{1}{2}x$}}
\shade{175}{2x-3y\leq6}{2x-3y\leq12}
\shade{176}{x+y\geq5}{-y+x\leq1}
\shade{177}{3x+7y\geq21}{x-y\leq2}
\shade{178}{4x-3y\leq12}{x\geq-\frac{3}{2}}
\shade{179}{3x+7y\geq21}{y\leq4}
\shade{180}{5x+7y\leq35}{x-2y\leq2}
\uq{181}{غیر مساواتوں کے نظام کا گراف بنا کر حل کا خطہ معلوم کریں: \lr{$x-2y\leq6,\ 2x+y\geq2$}}
\uq{182}{\lr{$z=3x+y$} کی زیادہ سے زیادہ اور کم سے کم قیمتیں معلوم کریں جب کہ: \lr{$3x+5y\geq15;\ x+6y\geq9;\ x\geq0;\ y\geq0$}}
\uq{183}{\lr{$f(x,y)=2x+5y$} کی زیادہ سے زیادہ قیمت معلوم کریں۔ جب کہ \lr{$2y-x\leq8,\ x-y\leq4;\ x\geq0;\ y\geq0$}}
\uq{184}{\lr{$f(x,y)=x+3y$} کی زیادہ سے زیادہ قیمت معلوم کریں جب کہ: \lr{$2x+5y\leq30,\ 5x+4y\leq20;\ x\geq0;\ y\geq0$}}
\uq{185}{\lr{$z=2x+3y$} کی زیادہ سے زیادہ قیمت معلوم کریں جب کہ: \lr{$2x+y\leq4;\ 4x-y\leq4;\ x\geq0;\ y\geq0$}}
\uq{186}{\lr{$z=2x+y$} کی کم سے کم قیمت معلوم کریں جب کہ: \lr{$x+y\geq3;\ 7x+5y\leq35;\ x\geq0;\ y\geq0$}}
\uq{187}{\lr{$f(x,y)=2x+3y$} کی زیادہ سے زیادہ قیمت معلوم کریں جب کہ: \lr{$2x+y\leq8;\ x+2y\leq14;\ x\geq0;\ y\geq0$}}
\uq{188}{درج ذیل غیر مساواتوں کے نظام کے لیے قابل عمل خطے کو سایہ دار کریں اور اس کے کونوں کے نقاط بھی معلوم کریں: \lr{$x-y\leq3,\ x+2y\leq6,\ x\geq0,\ y\geq0$}}
\uq{189}{تفاعل کی زیادہ سے زیادہ اور کم قیمتیں معلوم کریں جس کی وضاحت کی گئی ہے: \lr{$f(x,y)=2x+3y$} شرائط کے تابع \lr{$x-y\leq2,\ x+y\leq4,\ x\geq0,\ y\geq0.$}}
\uq{190}{حل کریں اور عددی خط پر ظاہر کریں: \lr{$5(x-3)\geq26x-(10x+4)$}}
\uq{191}{درج ذیل یک درجی غیر مساواتوں کے حل کا خطہ معلوم کریں: \lr{$3x-4y\leq12;\ 3x+2y\geq3$}}
\uq{192}{درج ذیل یک درجی غیر مساواتوں کے حل کا خطہ معلوم کریں: \lr{$2x+y\leq4;\ x+2y\leq6$}}
\uq{193}{\lr{$g(x,y)=x+4y$} کی زیادہ سے زیادہ قیمت معلوم کریں۔ جب کہ: \lr{$x+y\leq4,\ x\geq0$} اور \lr{$y\geq0$}.}
\uq{194}{\lr{$f(x,y)=3x+5y$} کی کم سے کم قیمت معلوم کریں۔ جب کہ \lr{$x+3y\geq3,\ x+y\geq2,\ x\geq0,\ y\geq0$}}
\uq{195}{اگر علاقہ کے مرکزی زاویے کی مقدار \lr{60$^{o}$} اور رداس \lr{12} سینٹی میٹر ہو تو اس علاقہ کا رقبہ معلوم کریں اور یہ دائرے کے کل رقبہ کا کتنا ہو گا؟}
\uq{196}{ایک دائروی علاقے کا رداس \lr{12} سینٹی میٹر ہے اور مرکزی زاویہ \lr{150$^{o}$} ہے۔ اس علاقے کو کاٹا گیا اور اس سے ایک کون بنائی گئی۔ اس کون کی ترچھی اونچائی اور قاعدے کا رداس معلوم کریں۔ اشارہ (علاقے کے قوس کی لمبائی = کون کا محیط)}
\trig{197}{\frac{\sin\theta}{\mathrm{cosec}\,\theta}+\frac{\cos\theta}{\sec\theta}=1}
\trig{198}{\frac{1-\sin\theta}{\cos\theta}=\frac{\cos\theta}{1+\sin\theta}}
\trig{199}{(\sec\theta-\tan\theta)^{2}=\frac{1-\sin\theta}{1+\sin\theta}}
\trig{200}{(\tan\theta+\cot\theta)^{2}=\sec^{2}\theta\,\mathrm{cosec}^{2}\theta}
\trig{201}{\frac{\tan\theta+\sec\theta-1}{\tan\theta-\sec\theta+1}=\tan\theta+\sec\theta}
\trig{202}{\sin^{3}\theta-\cos^{3}\theta=(\sin\theta-\cos\theta)(1+\sin\theta\cos\theta)}
\trig{203}{\sin^{6}\theta-\cos^{6}\theta=(\sin^{2}\theta-\cos^{2}\theta)(1-\sin^{2}\theta\cos^{2}\theta)}
\uq{204}{ثابت کریں:}\mline{\frac{1}{\mathrm{cosec}\,\theta-\cot\theta}-\frac{1}{\sin\theta}=\frac{1}{\sin\theta}-\frac{1}{\mathrm{cosec}\,\theta+\cot\theta}}
\uq{205}{ثابت کریں: \lr{$\sin^{6}\theta+\cos^{6}\theta=1-3\sin^{2}\theta\cos^{2}\theta$}}
\uq{206}{اگر \lr{$\tan\theta=\frac{3}{4}$} ہو تو باقی تکونیاتی نسبتیں معلوم کریں جب کہ \lr{$\theta$} پہلے ربع میں واقع ہو۔}
\uq{207}{پڑتال کریں: \lr{$\sin60^{o}\cos30^{o}+\cos60^{o}\sin30^{o}$}}
\quarter{208}{2\sin45^{o}-2\cos45^{o}}
\quarter{209}{3\cos45^{o}+4\sin45^{o}}
\quarter{210}{5\cos45^{o}-3\sin45^{o}}
\solvetri{211}{m\angle B=90^{\circ},\ m\angle C=60^{\circ},\ c=3\sqrt{3}\ cm}
\uq{212}{درج ذیل مثلث کو حل کریں جب کہ: \lr{$m\angle B=90^{\circ}$: a = 12 cm, C= 6 cm}}
\uq{213}{درج ذیل مثلث کو حل کریں جب کہ: \lr{$m\angle B=90^{\circ}$, b = 10cm, a = 6 cm}}
\uq{214}{سامنے دی گئی شکل میں \lr{x} کی لمبائی معلوم کریں}
\pfig{fig214}{5.5cm}
\uq{215}{درج ذیل مثلث کو حل کریں جب کہ: \lr{$m\angle B=90^{\circ},\ m\angle C=45^{\circ},$ a = 8 cm}}
\uq{216}{درج ذیل مثلث کو حل کریں جب کہ: \lr{$m\angle B=90^{\circ},\ m\angle A=60^{\circ},$ c = 4 cm}}
\uq{217}{درج ذیل مثلث کو حل کریں جب کہ: \lr{$m\angle B=90^{\circ},\ m\angle A=60^{\circ},$ c = 4 cm}}
\uq{218}{فرض کریں \lr{Q} اور \lr{R} دو نقاط نہر کے ایک ہی کنارے پر واقع ہے۔ ایک اور نقطہ \lr{P} نقطہ \lr{R} کی سیدھ میں دوسرے کنارے پر لیا گیا ہے۔ نہر کی چوڑائی معلوم کریں اور زاویہ \lr{PQR} کی مقدار ریڈین میں معلوم کریں۔}
\uq{219}{ایک سیڑھی دیوار کے ساتھ اس طرح کھڑی کی گئی ہے کہ اس کا پایا دیوار سے \lr{2} میٹر کے فاصلے پر ہے۔ اگر سیڑھی کی لمبائی \lr{8} میٹر ہو تو دیوار کی بلندی معلوم کریں۔}
\uq{220}{ایک مستطیلی کھیت کا وتر \lr{(x + 9)m} اور اضلاع \lr{(x + 7)m} اور \lr{xm} ہیں۔ \lr{x} کی قیمت معلوم کریں۔}
\pfig{fig220}{4.5cm}
\uq{221}{مثلث \lr{ABC} حل کریں جب کہ \lr{$m\angle B=90^{\circ},\ m\angle A=30^{\circ}$} اور \lr{a = 2}}
\uq{222}{مثلث \lr{ABC} حل کریں جب کہ \lr{$a=\sqrt{2}\ cm,\ $c = 1 cm} اور \lr{$m\angle B=90^{\circ}$}}
\uq{223}{ایک مساوی الساقین مثلث میں راسی زاویہ \lr{120$^{o}$} ہے۔ اگر مثلث کے قاعدہ کی لمبائی \lr{10 cm} ہو تو اس کے عمود کی مقدار معلوم کریں۔}
\uq{224}{سمندر کی سطح سے ایک روشنی کے مینار کی اونچائی \lr{150 m} ہے۔ اگر اس مینار کی اوپر والی سطح سے کسی بحری جہاز کا زاویہ صعود \lr{60$^{o}$} ہو تو اس مینار اور بحری جہاز کا درمیانی فاصلہ معلوم کریں۔}
\uq{225}{زمین کی سطح کے ایک نقطہ سے کسی ستون کی اوپر والی سطح کا زاویہ نزول \lr{15$^{o}$} ہے۔ اگر ستون کی طرف \lr{100m} چلنے کے بعد زاویہ نزول کی مقدار \lr{30$^{o}$} ہو جاتی ہے تو ستون کی اونچائی معلوم کریں۔}
\uq{226}{کسی پہاڑی کی چوٹی کا زاویہ نزول \lr{25$^{o}$} ہے۔ پہاڑی کی طرف \lr{100m} چلنے کے بعد زاویہ نزول کی مقدار \lr{45$^{o}$} ہو جاتی ہے۔ پہاڑی کی اونچائی معلوم کریں۔}
\uq{227}{\lr{300m} اونچی پہاڑی کی چوٹی سے نزدیکی دریا کے ساحل پر کسی نقطہ کا زاویہ صعود \lr{70$^{o}$} کا ہے۔ اور دریا کے دوسرے کنارے کا زاویہ صعود \lr{50$^{o}$} کا ہے۔ دریا کی چوڑائی معلوم کریں اور پہاڑی کی بنیاد سے دریا کا فاصلہ معلوم کریں۔}
\uq{228}{ایک پتنگ کے ساتھ \lr{120m} لمبا دھاگہ بندھا ہوا ہے۔ اور اس کا زاویہ نزول \lr{50$^{o}$} کا ہے۔ یہ پکڑنے والے کے ہاتھ سے کتنی دور ہے۔ (فرض کریں کہ دھاگہ تنا ہوا ہے)۔}
\uq{229}{ایک نگرانی کے مینار کے اوپر سے کسی عمارت کا زاویہ زمینی سطح پر \lr{45$^{o}$} ہے۔ اگر مینار کی اونچائی \lr{30m} ہو تو ایک آدمی مینار سے کتنا دور ہے؟}
\uq{230}{جب زمین کی سطح کے کسی نقطے سے \lr{40m} اونچے ستون کے اوپر والے سرے کو دیکھا جائے تو زاویہ نزول \lr{60$^{o}$} بنتا ہے۔ نقطے اور ستون کے پائے کا درمیانی فاصلہ معلوم کریں۔}
\trigb{231}{\sin\theta\,(\mathrm{cosec}\,\theta-\sin\theta)=\frac{1}{\sec^{2}\theta}}
\trigb{232}{\frac{\mathrm{cosec}\,\theta-\sec\theta}{\mathrm{cosec}\,\theta+\sec\theta}=\frac{1-\cos\theta}{1+\tan\theta}}
\trigb{233}{\frac{\cos\theta+\sin\theta}{\cos\theta-\sin\theta}+\frac{\cos\theta-\sin\theta}{\cos\theta+\sin\theta}=\frac{2}{1-2\sin^{2}\theta}}
\trigb{234}{\frac{1+\cos\theta}{1-\cos\theta}=(\mathrm{cosec}\,\theta+\cot\theta)^{2}}
\uq{235}{\lr{$\theta=\dfrac{3}{\sqrt{2}}$} ہو تو باقی تکونیاتی نسبتیں معلوم کریں جب \lr{$\theta$} پہلے ربط میں ہو۔}
\uq{236}{ثابت کریں: نقاط \lr{$A(0,2),\ B(\sqrt{3},1)$} اور \lr{$C(0,-2)$} قائمۃ الزاویہ مثلث کے راس ہیں۔}
\uq{237}{ثابت کریں کہ: \lr{$A(3,1),\ B(-2,-3)$} اور \lr{$C(2,2)$} متساوی الساقین مثلث کے راس ہیں۔}
\uq{238}{ثابت کریں کہ: \lr{$A(5,2),\ B(-2,3),\ C(-3,-4)$} اور \lr{$D(4,-5)$} متوازی الاضلاع کے راس ہیں۔}
\uq{239}{\lr{h} کی قیمت معلوم کریں جب کہ نقاط \lr{$B(0,2),\ A(\sqrt{3},-1)$} اور \lr{$C(h,-2)$} قائمۃ الزاویہ مثلث کے راس ہیں جس میں قائمۃ الزاویہ راس \lr{A} پر ہے۔}
\uq{240}{نقاط، \lr{$A(-1,h),\ B(3,2)$} اور \lr{$C(7,3)$} ہم خط ہیں۔ \lr{h} معلوم کریں۔}
\uq{241}{نقاط \lr{$A(-5,-2)$} اور \lr{$B(5,-4)$} دائرے کے قطر کے سرے ہیں۔ دائرے کا مرکز اور رداس معلوم کریں۔}
\uq{242}{نقاط \lr{$A(h,1),\ B(2,7)$} اور \lr{$C(-6,-7)$} قائمۃ الزاویہ مثلث کے راس ہیں جس میں قائمۃ الزاویہ راس \lr{A} پر ہے۔ \lr{h} معلوم کریں۔}
\uq{243}{نقاط \lr{$A(9,3),\ B(-7,7),\ C(-3,-7)$} اور \lr{$D(5,-5)$} چوکور کے راس ہیں۔ اس کے اضلاع کے وسطی نقاط معلوم کریں، ثابت کیجیے کہ وسطی نقاط کو ملانے سے جو شکل بنتی ہے وہ متوازی الاضلاع ہے۔}
\uq{244}{دیے گئے نقاط کے درمیان فاصلہ معلوم کریں: \lr{$C(-4,-2),\ D(0,9)$}}
\uq{245}{ڈھلوانوں کا استعمال کرتے ہوئے ثابت کریں کہ مثلث اپنے راس \lr{$A(6,1),\ B(2,7)$} اور \lr{$C(-6,-7)$} کے ساتھ ایک قائمۃ الزاویہ مثلث بناتی ہے۔}
\ptsq{246}{(1,-2),\ (2,4)\text{ اور }(4,1),\ (-8,2)}
\ptsq{247}{(-3,4),\ (6,2)\text{ اور }(4,5),\ (-2,-7)}
\uq{248}{نقاط \lr{$A(3,5)$} اور \lr{$B(9,8)$} کو ملانے والے قطعہ خط کے عمودی ناصف کی مساوات معلوم کریں۔}
\uq{249}{\lr{$(-4,-6)$} سے گزرنے والے خط کی مساوات معلوم کریں جو کہ اس خط پر عمود ہو اور جس کی ڈھلوان \lr{$\frac{-3}{2}$} ہو۔}
\uq{250}{\lr{$(11,-5)$} سے گزرنے والے خط کی مساوات معلوم کریں جو کہ اس خط کے متوازی ہو جس کی ڈھلوان \lr{-24} ہو۔}
\lineform{251}{2x-4y+11=0}
\lineform{252}{4x+7y-2=0}
\lineform{253}{15y-8x+3=0}
\parperp{254}{2x+y-3=0;\ 4x+2y+5=0}
\parperp{255}{3y=2x+5;\ 3x+2y-8=0}
\parperp{256}{4y+2x-1=0;\ x-2y-7=0}
\uq{257}{\lr{$(-4,7)$} سے گزرنے والے خط کی مساوات معلوم کریں جو کہ خط \lr{$2x-7y+4=0$} کے متوازی ہو}
\uq{258}{\lr{$(5,-8)$} سے گزرنے والے خط کی مساوات معلوم کریں جو \lr{$A(-15,-8),\ B(10,7)$} سے بننے والے خط پر عمود ہو۔}
\uq{259}{خط مستقیم کی مساوات معلوم کریں اگر اس کی ڈھلوان \lr{2} اور \lr{y}-قاطع \lr{5} ہو۔}
\uq{260}{خط مستقیم کی مساوات معلوم کریں اگر: یہ ڈھلوان \lr{-6} والے خط پر عمود ہو اور \lr{y}-قاطع \lr{$\frac{4}{3}$} ہو۔}
\uq{261}{نقطہ \lr{(5,1)} سے گزرنے والے خط مستقیم کی ایک مساوات لکھیں جو نقاط \lr{(0,1),(7,15)} سے گزرنے والے خط کے متوازی ہو۔}
\uq{262}{مبدا سے ایک خط تک کی عمودی لمبائی \lr{5} سینٹی میٹر ہے۔ اور اس عمود کا جھکاؤ \lr{120$^{o}$} ڈگری ہے۔ خط کی ڈھلوان اور \lr{y}-قاطع معلوم کریں۔}
\uq{263}{ایک معمار ایک پارک ڈیزائن کر رہا ہے جس میں دو عمارتوں کے محدودات \lr{(10, 8)} اور \lr{(4, 3)} گرڈ پر واقع ہیں۔ عمارتوں کے درمیان فاصلہ میٹروں میں معلوم کریں۔}
\uq{264}{ڈیلیوری ڈرائیور کو دو ڈیلیوری مقامات کے درمیان فاصلہ معلوم کرنے کی ضرورت ہے۔ ایک مقام شہر کے گرڈ کے نقشے پر \lr{(7, 2)} اور دوسرا \lr{(12, 10)} پر واقع ہے۔ دونوں مقامات کے درمیان فاصلہ کلو میٹر میں معلوم کریں۔}
\uq{265}{ایک جہاز \lr{$(12^{\circ}N,\ 65^{\circ}W)$} پر واقع پورٹ \lr{A} سے پورٹ \lr{B} تک \lr{$(20^{\circ}N,\ 45^{\circ}W)$} جا رہا ہے۔ اگر جہاز زمین کی سطح پر مختصر ترین راستے پر سفر کرتا ہے تو نقاط کے درمیان خط مستقیم کا فاصلہ معلوم کریں۔}
\uq{266}{فرح ایک مستطیل میدان کے گرد باڑ لگا رہی ہے جس کے کونوں کے محدودات \lr{(0, 0), (0, 5), (8, 5)} اور \lr{(8, 0)} ہیں۔ اس کھیت کے احاطے کے لیے کتنی باڑ کی ضرورت ہو گی؟}
\uq{267}{ایک ہوائی جہاز شہر \lr{X} \lr{$(40^{\circ}N,\ 100^{\circ}W)$} سے شہر \lr{Y} \lr{$(50^{\circ}N,\ 80^{\circ}W)$} تک اڑ رہا ہے۔ محدداتی جیومیٹری کا استعمال کرتے ہوئے ان دو شہروں کے درمیان فاصلہ معلوم کریں۔}
\uq{268}{زمین کا سروے کرنے والا ایک مستطیلی پلاٹ کو نشان زدہ کر رہا ہے جس کے کونوں کے محدودات \lr{(3, 1), (3, 6), (8, 6)} اور \lr{(8, 1)} ہیں۔ اس مستطیلی پلاٹ کا احاطہ معلوم کریں۔}
\uq{269}{ایک مستطیلی باغ جس کے کونوں کے محدودات \lr{$A(0,0),\ B(5,0),\ C(5,3)$} اور \lr{$D(0,3)$} ہیں۔ اس مستطیلی باغ کے گرد کتنی باڑ لگانے کی ضرورت ہے؟}
\uq{270}{زمین کی تزئین و آرائش کرنے والا ایک مثلثی باغ ڈیزائن کر رہا ہے جس کے کونے \lr{A(2,3), B(5,7)} اور \lr{C(6,2)} پر ہیں۔ مثلث کے اضلاع کی لمبائیاں معلوم کریں۔}
\uq{271}{نقاط \lr{(3,7)} اور \lr{(5,11)} سے گزرنے والے خط کی مساوات \lr{y = mx + c} شکل میں معلوم کریں۔}
\uq{272}{ایک ڈرون ایک مقام سے دوسرے مقام تک اڑ رہا ہے جس کے محدودات \lr{(2,3)} اور \lr{(10,15)} ہیں۔ اس خط کی ڈھلوان اور کل فاصلہ معلوم کریں۔}
\dproof{273}{(x-4)^{2}+9=x^{2}-8x+25}
\dproof{274}{(x+1)^{2}-(x-1)^{2}=4x}
\dproof{275}{(x+5)^{2}-(x-5)^{2}=20x}
\jproof{276}{\frac{4+16x}{4}=1+4x}
\jproof{277}{\frac{6x^{2}+18x}{3x^{2}-9}=\frac{2x}{x-3}}
\jproof{278}{\frac{x^{2}+7x+10}{x^{2}-3x-9}=\frac{x+5}{x-5}}
\uq{279}{درج ذیل بیان کو ثابت کریں۔ اگر \lr{x} ایک طاق صحیح عدد ہو تو ثابت کریں کہ \lr{$x^{2}-4x+6$} ایک طاق عدد ہے۔}
\uq{280}{درج ذیل بیان کو ثابت کریں۔ اگر \lr{x} ایک طاق صحیح عدد ہو تو ثابت کریں کہ \lr{$x^{2}+2x+4$} ایک جفت عدد ہے۔}
\uq{281}{ثابت کریں کہ کوئی سے بھی دو غیر خالی سیٹوں \lr{A} اور \lr{B} کے لیے، \lr{$(A\cap B)'=A'\cup B'$}}
\uq{282}{ثابت کریں کہ \lr{$[(p\rightarrow q)\wedge p]$} اور \lr{$[(p\rightarrow q)\wedge p]\rightarrow q$} کا صدق جدول بنائیں۔}
\uq{283}{ثابت کریں کہ کوئی سے دو خالی سیٹوں \lr{A} اور \lr{B} کے لیے \lr{$(A\cup B)'=A'\cap B'$}}
\uq{284}{ثابت کریں کہ \lr{$\dfrac{a}{b}+\dfrac{c}{d}=\dfrac{ad+bc}{bd}$} جبکہ \lr{a,b,c} اور \lr{d} غیر صفر حقیقی اعداد ہیں۔}
\uq{285}{ثابت کریں: \lr{$(x+1)^{2}+7=x^{2}+2x+8$}}
\uq{286}{ہر مرحلے کا جواز پیش کر کے ثابت کریں کہ \lr{$\dfrac{45x+15}{15}=3x+1$}}
\uq{287}{لکڑی کا ایک تختہ سیدھا سیڑھی کے اوپر رکھا گیا ہے جس کی چوڑائی \lr{20} سینٹی میٹر اور گہرائی \lr{16} سینٹی میٹر ہے۔ ایک مستطیلی صندوق جس کی اونچائی \lr{8} سینٹی میٹر اور چوڑائی \lr{x} سینٹی میٹر ہے۔ تختے کے نیچے سیڑھی پر رکھا گیا ہے۔ \lr{x} کی قیمت معلوم کریں۔}
\uq{288}{ایک آدمی جس کا قد \lr{1.8} میٹر ہے۔ اس کے سائے کی لمبائی \lr{0.76} میٹر ہے۔ اگر عین اسی وقت ایک ٹیلی فون کے کھمبے کا سایہ \lr{3} میٹر ہو تو کھمبے کی اونچائی معلوم کریں۔}
\uq{289}{ایک متساوی الساقین ذونقہ \lr{ABCD} بنائیے۔ جس میں \lr{$\overline{AB}\parallel\overline{CD}$} اور \lr{$m\overline{AB}>m\overline{CD}$}، وتر \lr{$\overline{AC}$} اور \lr{$\overline{BD}$} اس طرح کھینچیں کہ وہ ایک دوسرے کو نقطہ \lr{E} پر قطع کرتے ہوں۔ ثابت کیجیے \lr{$\triangle ABE$} اور \lr{$\triangle CDE$} متشابہ ہیں۔ اگر \lr{$m\overline{AB}=8\ cm$}، \lr{$m\overline{CD}=4\ cm$} اور \lr{$m\overline{AE}=3\ cm$} ہو تو \lr{$\overline{CE}$} کی لمبائی معلوم کریں۔}
\uq{290}{ایک منظم بارہ اضلاع والی شکل کی لمبائیاں \lr{$\frac{1}{\sqrt{2}}$} کے جزو ضربی کی نسبت سے کم کی گئی ہیں۔ اگر اصل منظم بارہ اضلاع والی شکل کا احاطہ \lr{72cm} ہو تو چھوٹی بارہ اضلاع والی شکل کے ضلع کی مقدار معلوم کریں۔}
\uq{291}{سامنے دی گئی شکل میں \lr{$\triangle ABC\sim\triangle DEF$} میں \lr{$m\overline{AB}=12cm,\ m\overline{AC}=20cm,\ m\overline{BC}=16cm$}۔ \lr{$\triangle DEF$} میں \lr{$m\overline{DE}=6cm$} ہے۔ \lr{$m\overline{DE}$} اور \lr{$m\overline{EF}$} کی لمبائیاں معلوم کریں۔}
\pfig{fig291}{5.5cm}
\uq{292}{دو متوازی الاضلاع میں دو اضلاع کا درمیانی زاویہ \lr{45$^{o}$} ہو تو معلوم کریں کہ کیا دونوں متوازی الاضلاع متشابہ ہیں؟}
\uq{293}{مثلثان \lr{XBC} اور \lr{XDE} میں \lr{x} اور \lr{y} کی مقداریں معلوم کریں۔}
\uq{294}{چوکور \lr{ABCD} کے اضلاع کی مقداریں \lr{$m\overline{AB}=5cm,\ m\overline{BC}=8,\ m\overline{CD}=10\ cm,\ m\overline{AD}=12\ cm$} اور اس کے زاویے \lr{$m\angle A=90^{\circ},\ m\angle B=120^{\circ}$} اور \lr{$m\angle C=90^{\circ}$} ہیں اور چوکور \lr{EFGH} کے اضلاع کی مقداریں \lr{$m\overline{EF}=10cm,\ m\overline{FG}=16cm,\ m\overline{GH}=20cm,\ m\overline{EH}=24cm$} اور اس کے زاویے \lr{$m\angle E=90^{\circ},\ m\angle F=120^{\circ}$} اور \lr{$m\angle H=60^{\circ}$} ہیں۔ تو ثابت کریں کہ چوکور \lr{ABCD} اور \lr{EFGH} متشابہ ہیں۔}
\uq{295}{\lr{$\triangle ABC=36\ cm^{2}$} ہے اور \lr{$m\overline{AB}=6\ cm,\ m\overline{BD}=4\ cm$}، \lr{$\triangle ADE$} کا رقبہ معلوم کریں۔}
\uq{296}{\lr{$\triangle ABC=36\ cm^{2}$} ہے اور \lr{$m\overline{AB}=6\ cm,\ m\overline{BD}=4\ cm$}، ذونقہ \lr{BCED} کا رقبہ معلوم کریں۔}
\uq{297}{\lr{$\triangle ABC$} اور \lr{$\triangle DEF$} سکیل \lr{k=3} متناسبی عامل کے ساتھ متشابہ ہیں اگر \lr{$\triangle ABC$} کا رقبہ \lr{50 $cm^{2}$} ہو تو \lr{$\triangle DEF$} کا رقبہ معلوم کریں۔}
\uq{298}{چوکور \lr{ABCD} اور چوکور \lr{EFGH} سکیل فیکٹر \lr{$k=\frac{1}{4}$} کے ساتھ متشابہ ہیں۔ اگر چوکور \lr{ABCD} کا رقبہ \lr{64 $cm^{2}$} ہو تو چوکور \lr{EFGH} کا رقبہ معلوم کریں۔}
\uq{299}{دو متشابہ مثلثوں کے رقبے \lr{16 $cm^{2}$} اور \lr{25 $cm^{2}$} ہیں۔ متناظرہ اضلاع کے جوڑے کی نسبت معلوم کریں۔}
\uq{300}{دو متشابہ مثلثوں کے رقبے \lr{144 $cm^{2}$} اور \lr{81 $cm^{2}$} ہیں اگر بڑی مثلث کے قاعدہ کی لمبائی \lr{30cm} ہو تو چھوٹی مثلث کے متناظرہ قاعدہ کی لمبائی معلوم کریں۔}
\uq{301}{ایک منظم سات اضلاع والی شکل ایک بڑی منظم سات اضلاع والی شکل میں محصور ہے اور بڑی سات اضلاع والی شکل کے ہر ایک ضلع کی لمبائی چھوٹی سات اضلاع والی شکل کے ہر ضلع کا \lr{1.7} گنا ہے۔ اگر چھوٹی سات اضلاع والی شکل کا رقبہ \lr{100 $cm^{2}$} ہو تو بڑی سات اضلاع والی شکل کا رقبہ معلوم کریں۔}
\uq{302}{دو متشابہ سلنڈر نما پانی کے تالابوں کی اونچائیوں میں نسبت \lr{5:3} ہے۔}
\usc{\lr{1.} اگر بڑے تالاب کی سطح کا رقبہ \lr{250} سکویر میٹر ہو تو چھوٹے تالاب کی سطح کا رقبہ معلوم کریں۔}
\usc{\lr{2.} اگر چھوٹے تالاب کا حجم \lr{270} کیوبک میٹر ہو تو بڑے تالاب کا حجم معلوم کریں۔}
\uq{303}{دو متشابہ مخروطی کینز کی متناظرہ لمبائیوں میں نسبت \lr{3:2} ہے۔}
\usc{\lr{(i)} بڑی مخروطی کون کی سطح کا رقبہ \lr{96 m$^{2}$} ہو تو چھوٹی مخروطی کین کی سطح کا رقبہ معلوم کریں۔}
\usc{\lr{(ii)} چھوٹی مخروطی کین کا حجم \lr{240 m$^{3}$} ہو تو بڑی مخروطی کین کا حجم معلوم کریں۔}
\uq{304}{دو متشابہ سلنڈر نما کینز کی متناظرہ لمبائیوں میں نسبت \lr{3:2} ہے۔ بڑے سلنڈر نما کین کی سطح کا رقبہ \lr{67.5} سکویر میٹر ہے۔ چھوٹے سلنڈر نما کین کی سطح کا رقبہ معلوم کریں۔}
\uq{305}{دو متشابہ سلنڈر نما کین کی متناظرہ لمبائیوں میں نسبت \lr{3:2} ہے۔ چھوٹے سلنڈر نما کین کا حجم \lr{132} کیوبک میٹر ہے۔ چھوٹے سلنڈر نما کین کا حجم معلوم کریں۔}
\uq{306}{ایک مجسم کون \lr{C} کو دو حصوں \lr{A} اور \lr{B} میں اس طرح کاٹا گیا ہے کہ اس کے ڈھلوانی کنارے \lr{6} سینٹی میٹر اور \lr{4} سینٹی میٹر ہیں۔ نسبت معلوم کریں: کون \lr{A} اور کون \lr{C} کے قاعدے کے قطروں میں}
\uq{307}{ایک مجسم کون \lr{C} کو دو حصوں \lr{A} اور \lr{B} میں اس طرح کاٹا گیا ہے کہ اس کے ڈھلوانی کنارے \lr{6} سینٹی میٹر اور \lr{4} سینٹی میٹر ہیں۔ کون \lr{A} اور کون \lr{C} کے حجم میں}
\uq{308}{ایک مجسم کون \lr{C} کو دو حصوں \lr{A} اور \lr{B} میں اس طرح کاٹا گیا ہے کہ اس کے ڈھلوانی کنارے \lr{6} سینٹی میٹر اور \lr{4} سینٹی میٹر ہیں۔ کون \lr{A} اور کون \lr{C} کے قاعدے کے رقبوں میں}
\uq{309}{ایک مجسم کون \lr{C} کو دو حصوں \lr{A} اور \lr{B} میں اس طرح کاٹا گیا ہے کہ اس کے ڈھلوانی کنارے \lr{6} سینٹی میٹر اور \lr{4} سینٹی میٹر ہیں۔ اگر کون \lr{A} کا حجم \lr{72 cm$^{3}$} ہو تو کون \lr{B} کا حجم معلوم کریں۔}
\uq{310}{چاولوں کی ایک بوری کا وزن \lr{50} کلو گرام اور اونچائی \lr{5} میٹر ہے۔ اسی طرح کی بوری کا ماس معلوم کریں جس کی اونچائی \lr{6} میٹر ہے۔}
\uq{311}{ایک مربع شکل کو وتر کی طرف دو آدھے حصوں میں کاٹ کر قائمۃ الزاویہ مثلثوں میں سے ہر ایک کے وتر کی دوسری کے وتر سے جوڑ کر ایک نئی شکل بنائی گئی وضاحت کریں یہ شکل کیوں اقلیدسی نمونہ ہے۔ اور اس نئی شکل کا زاویہ معلوم کریں۔}
\uq{312}{بنیادی اشکال کو بار جوڑ کر ایک اقلیدسی نمونہ بنایا گیا ہے بنیادی شکل ایک قائمۃ الزاویہ مثلث ہے جس کے اضلاع کی مقدار \lr{3،4} اور \lr{5} اکائیاں ہیں۔ ایک مربع شکل جس کا رقبہ \lr{3600} مربع اکائیاں ہوں اس پر اقلیدسی نمونہ بنانے کے لیے قائمۃ الزاویہ مثلث کے کم سے کم کتنے عکس درکار ہوں گے؟}
\uq{313}{منظم مسدسوں کو استعمال کرکے ایک اقلیدسی نمونہ بنایا گیا ہے۔ ہر ایک مسدس کئے ضلع کی لمبائی \lr{5 cm} ہے۔ افر اس اقلیدسی نمونہ \lr{25} مسدسوں پر مشتمل ہو تو اس کا رقبہ معلوم کریں اور فرض کریں یہ اقلیدسی مکمل مسدس ہے تو اس کے بیرونی کنارے کا احاطہ معلوم کریں۔}
\uq{314}{ایک متوازی الاضلاع \lr{ABCD} میں اگر \lr{$m\angle DAB=70^{\circ}$} ہو تو متوازی الاضلاع میں باقی زاویوں کی مقداریں معلوم کریں۔}
\uq{315}{ایک ذونقہ شکل کی کھڑکی جس کے متوازی اضلاع کی لمبائیاں \lr{3} میٹر اور \lr{1.5} میٹر ہیں اور اونچائی \lr{2} میٹر ہے۔ کھڑکی کا رقبہ معلوم کریں۔}
\uq{316}{فلک \lr{100} سکویر میٹر کے مربعی شکل کے صحن میں ڈائل لگوانے کا منصوبہ بناتا ہے۔ وہ \lr{0.25} سکویر میٹر کی مربع شکل اور مثلثی شکل دو طرح کی ٹائل لگوانے کا فیصلہ کرتا ہے۔ اگر مربع شکل کی ٹائلیں \lr{60} فی صد اور مثلثی شکل کی \lr{40} فی صد ہوں تو ہر ایک شکل کی کتنی ٹائلوں کی ضرورت ہوگی؟}
\uq{317}{علم منظم مسدس (ہر ایک ضلع کی لمبائی \lr{1} میٹر) اور مساوی الاضلاع مثلث (ہر ایک ضلع کی لمبائی \lr{1} میٹر) استعمال کرتے ہوئے مستطیلی رقبہ \lr{5} میٹر بائے \lr{10} میٹر پر فرشی ڈیزائن بنانا چاہتا ہے۔ اس رقبہ پر اقلیدسی نمونہ بنانے کے لیے علی کے کتنے مسدسوں اور مثلثوں کی ضرورت ہوگی؟}
\uq{318}{ایک متوازی الاضلاع شکل کا کمرہ جس کے قاعدے کی لمبائی \lr{10} میٹر اور اونچائی \lr{8} میٹر ہے۔ بابر قالین کا رول استعمال کرتے ہوئے کمرے کو کارپٹ کرنا چاہتا ہے۔ ایک رول \lr{20 m$^{2}$} جگہ گھیر تا ہے۔ اس کمرہ کو کارپٹ کرنے کے لیے قالین کے کتنے رول درکار ہوں گے؟}
\uq{319}{پانچ اضلاع والی اور دس اضلاع والی اشکال کے مجموعہ کو استعمال کرتے ہوئے ایک اقلیدسی شکل بنائی گئی ہے۔ اس کے راس کے زاویوں کا مجموعہ معلوم کریں جہاں پر پانچ اضلاع والی اور دس اضلاع والی اشکال ملتی ہیں۔}
\uq{320}{درج ذیل دو درجی اور سہ درجی تفاعل کا گراف بنائیں:}\mline{y=x^{3}+2x^{2}-5x-6;\ -3.5\leq x\leq2.5}
\plotq{321}{y=x^{2}+x-2}
\plotq{322}{y=x^{3}+3x^{2}+2x;\ -2.5\leq x\leq0.5}
\plotq{323}{y=5x^{2}-2x-3}
\uq{324}{\lr{$y=2x^{2}-4x+3$} کے گراف کو \lr{x} کی قدروں \lr{-1} سے \lr{3} تک بنائیں۔ \lr{(2,3)} پر مماس بنائیں اور ڈھلوان معلوم کریں۔}
\uq{325}{\lr{$y=3x^{2}+x+1$} کا گراف بنائیں اور \lr{(1,5)} پر مماس کھینچیں۔ اس نقطہ پر مماس کی ڈھلوان بھی معلوم کریں۔}
\uq{326}{\lr{2016} میں ایک اسکول میں طلبہ کی تعداد \lr{1000} تھی۔ اگر مساوات \lr{$S=1000\,e^{-1}$} کے مطابق تعداد زوال پذیر ہوتی ہے، جہاں \lr{t} وقت پر طلبہ کی تعداد \lr{S} ہوتی ہے۔ \lr{t=0 (in 2016)} سے \lr{t=9 (in 2025)} کے لیے دی گئی مساوات کا گراف بنائیں۔}
\uq{327}{\lr{2016} میں ایک اسکول میں طلبہ کی تعداد \lr{1000} تھی۔ اگر مساوات \lr{$S=1000\,e^{-1}$} کے مطابق تعداد زوال پذیر ہوتی ہے، جہاں \lr{t} وقت پر طلبہ کی تعداد \lr{S} ہوتی ہے۔ گراف سے \lr{2019} اور \lr{2023} میں طلبہ کی تعداد معلوم کریں۔}
\uq{328}{ایک اخباری ایجنسی نے \lr{70} روپے فی ایڈیشن اور معمولی پرنٹنگ اور تقسیم کے اخراجات \lr{40} روپے فی کاپی کی قیمت مقرر کی ہے۔ منافع کا تفاعل \lr{$p(x)=0.10x+70$} ہے۔ جہاں \lr{x} اخبارات کی تعداد ہے۔ گراف بنائیں اور \lr{500} اخبارات کے لیے منافع معلوم کریں۔}
\uq{329}{علی اسکول کو فروخت کے لیے مہنگی قمیضیں تیار کرتا ہے۔ \lr{x} قمیضوں کی قیمت (روپوں) میں \lr{$C(x)=1500+10x+0.2x^{2},\ 0\leq x\leq150$} ہے۔ گراف بنائیں اور \lr{200} قمیضوں کی قیمت معلوم کریں۔}
\uq{330}{\lr{2010} میں ایک گاؤں کی آبادی \lr{753} تھی۔ اگر مساوات \lr{$p=753e^{0.03t}$} کے مطابق آبادی بڑھتی ہے، جہاں \lr{p} آبادی میں \lr{t} وقت میں افراد کی تعداد ہے۔ \lr{t=0 (in 2010)} سے \lr{t = 30 (in 2040)} کے لیے آبادی کی مساوات کا گراف بنائیں۔}
\uq{331}{\lr{2010} میں ایک گاؤں کی آبادی \lr{753} تھی۔ اگر مساوات \lr{$p=753e^{0.03t}$} کے مطابق آبادی بڑھتی ہے، جہاں \lr{p} آبادی میں \lr{t} وقت میں افراد کی تعداد ہے۔ گراف سے آبادی کا تخمینہ لگائیں (الف) \lr{2020} میں اور (ب) \lr{2030} میں۔}
\uq{332}{\lr{$y=x^{2}-3x-2$} کے گراف کو \lr{x} کی قیمتوں کے لیے \lr{-8} سے \lr{8} تک بنائیں \lr{x = -6} پر مماس کھینچیں اور ڈھلوان کا تعین کریں۔}
\uq{333}{ماجد کی تنخواہ \lr{S(x)} روپوں میں درج ذیل کلیے پر مبنی ہے: \lr{$S(x)=25000+1500x$} جہاں \lr{x} ان سالوں کی تعداد ہے جو اس نے کام کیا \lr{$10\leq x\leq0$} کے لیے تنخواہ کے تفاعل کا گراف بنائیں اور اس کی تشریح کریں۔}
\uq{334}{ایک کمپنی فٹ بال تیار کرتی ہے۔ \lr{x} فٹ بال تیاری کی لاگت \lr{$C(x)=90,000+600x$} ہے۔ \lr{x} فٹ بالوں کی فروخت سے حاصل ہونے والی آمدنی \lr{$R(x)=1,800x$} ہے۔ مساوی نقطہ معلوم کریں اور \lr{200} فٹ بال فروخت ہونے پر منافع یا نقصان کا تعین کریں۔ دونوں تفاعل کے گراف کھینچیں اور مساوی نقطہ کی نشاندہی کریں۔}
\uq{335}{ایک نئے رسالے کی فروخت کے بڑھنے کی توقع کو مساوات کے ذریعے دکھایا گیا ہے: \lr{$S=200000\left(1-e^{-0.05t}\right)$} یہاں \lr{t} ہفتوں میں دیا گیا ہے۔ پہلے \lr{50} ہفتوں کے لیے فروخت کا گراف بنائیں۔}
\uq{336}{ایک نئے رسالے کی فروخت کے بڑھنے کی توقع کو مساوات کے ذریعے دکھایا گیا ہے: \lr{$S=200000\left(1-e^{-0.05t}\right)$} یہاں \lr{t} ہفتوں میں دیا گیا ہے۔ فروخت ہونے والے رسالوں کی تعداد معلوم کریں، جب \lr{t=5} اور \lr{t=35} ہو۔}
\uq{337}{\lr{x} کی \lr{-5} سے \lr{4} تک کے لیے \lr{$y=\dfrac{1}{2}(x+4)(x-1)(x-3)$} کا گراف بنائیں۔}
\uq{338}{کسی خاص مارکیٹ کے لیے طلب اور رسد کے تفاعل مساواتوں کے ذریعے دیے گئے ہیں: \lr{$P_{s}=Q^{2}+5$} اور \lr{$P_{d}=Q^{2}-10Q$} جہاں \lr{P} قیمت اور \lr{Q} تعداد کو ظاہر کرتا ہے۔ وقعہ \lr{Q=-20} سے \lr{Q=20} پر ہر تفاعل کا گراف بنائیں۔}
\uq{339}{ایک ٹیلی ویژن بنانے والی کمپنی \lr{40} انچ ایل ای ڈی بناتی ہے۔ \lr{x} ایل ای ڈیز کی تیاری کی لاگت \lr{$C(x)=60,000+250x$} ہے اور \lr{x} ایل ای ڈیز کی فروخت سے حاصل ہونے والی آمدنی \lr{$R(x)=1200x$} ہے۔ مساوی نقطہ معلوم کریں اور \lr{100} ایل ای ڈیز فروخت ہونے پر نفع یا نقصان معلوم کریں۔ مساوی نقطہ گراف پر ظاہر کریں۔}
\bisect{340}{m\overline{AC}=5.3\ cm\ \text{اور}\ m\overline{AB}=4.5\ cm,\ m\angle A=45^{\circ}}
\bisect{341}{m\angle B=60^{\circ}\ \text{اور}\ m\overline{AB}=6cm,\ m\angle A=150^{\circ}}
\perpbis{342}{m\overline{AC}=7\ cm\ \text{اور}\ m\overline{AB}=5\ cm,\ m\overline{BC}=6cm}
\perpbis{343}{m\overline{BC}=6.5\ cm\ \text{اور}\ m\overline{AB}=7.1\ cm,\ m\angle B=135^{\circ}}
\medq{344}{m\angle M=38^{\circ}\ \text{اور}\ m\overline{LM}=4.9\ cm,\ m\angle L=51^{\circ}}
\medq{345}{m\overline{LM}=8.1cm\ \text{اور}\ m\overline{MN}=4.8\ cm,\ m\angle N=30^{\circ}}
\bisect{346}{m\overline{AC}=5.3\ cm\ \text{اور}\ m\overline{AB}=4.5\ cm,\ m\angle A=45^{\circ}}
\uq{347}{\lr{$\triangle DEF$} کی پیمائشیں دی گئی ہیں۔}\mline{m\angle E=45^{\circ}\ \text{اور}\ \triangle DEF:\ m\overline{DE}=4.8\ cm,\ m\overline{EF}=4\ cm}
\usc{\lr{$\triangle DEF$} کی عمود کھنچیں اور عمودی مرکز معلوم کریں۔}
\uq{348}{درج ذیل مثلث بنائیں اور معلوم کریں کہ کیا کوئی مبہم صورت موجود ہے؟}\mline{m\overline{CD}=4.7\ cm\ \text{اور}\ \triangle BCD;\ m\overline{BC}=5\ cm,\ m\angle B=62^{\circ}}
\uq{349}{درج ذیل مثلث بنائیں اور معلوم کریں کہ کیا کوئی مبہم صورت موجود ہے؟}\mline{m\overline{LN}=5\ cm\ \text{اور}\ \triangle KLM;\ m\overline{LM}=6\ cm,\ m\angle M=42^{\circ}}
\uq{350}{ایک مثلث بنائیں جس کے اضلاع کی پیمائش \lr{5.3 cm, 5.9 cm} اور \lr{6.2 cm} ہو۔}
\uq{351}{ایک مثلث \lr{BCD} بنائیں جس میں دو اضلاع کی لمبائیاں \lr{5.5 cm} اور \lr{4.2} سینٹی میٹر ہوں اور ان کا درمیانی زاویہ \lr{60$^{o}$} ہو۔}
\uq{352}{ایک مثلث \lr{CDE} بنائیں جب \lr{$m\overline{DE}=4.3\ cm,\ m\angle D=30^{\circ}$} اور \lr{$m\angle E=120^{\circ}$} ہو۔}
\uq{353}{ایک مثلث \lr{DEF} بنائیں جب \lr{$m\overline{DE}=6\ cm,\ m\angle D=110^{\circ}$} اور \lr{$m\overline{EF}=9\ cm$} ہو۔}
\uq{354}{مثلثان \lr{DEF} اور \lr{$DEF'$} بنائیں جب \lr{$m\overline{DE}=6\ cm,\ m\angle D=30^{\circ}$} اور \lr{$m\overline{EF}=3.6\ cm$} ہو۔}
\uq{355}{\lr{null}}
\uq{356}{\lr{null}}
\uq{357}{مثلث \lr{EFG} کے عمودی ناصف اور ان کے وسطانیے کھینچیں۔ جس میں \lr{$m\overline{EG}=4.3\ cm,\ m\overline{EF}=5\ cm,\ m\overline{FG}=2.5\ cm$} ہو۔}
\uq{358}{مثلث \lr{FGH} کے زاویوں کے ناصف کھینچیں اگر \lr{$m\overline{FG}=5.2,\ m\overline{GH}=4.1$} اور \lr{$m\angle FGH=120^{\circ}$} ہو۔}
\uq{359}{مثلث \lr{GHI} بنائیں جس میں \lr{$m\angle G=50^{\circ},\ m\overline{GH}=5.7,\ m\angle G=68^{\circ}$} ہو۔ ثابت کریں کہ مثلث \lr{GHI} کے ارتفاع کا ہم نقطہ ہوتے ہیں۔}
\uq{360}{ایک مثلث \lr{CDE} بنائیں۔ جس میں \lr{$m\overline{CD}=7.6\ cm,\ m\angle D=45^{\circ}$} اور \lr{$m\overline{DE}=5.9\ cm$} ہو۔ تمام نقاط کا لوسائی کھنچیں جو کہ \lr{C} اور \lr{D} سے برابر فاصلے پر ہوں۔ نقطہ \lr{X} کی نشاندہی کریں جہاں یہ دونوں لوسائی ایک دوسرے کو قطع کرتے ہیں}
\uq{361}{دو نقاط \lr{A} اور \lr{B} کے درمیان فاصلہ \lr{8.2} سینٹی میٹر ہے۔ نقطہ \lr{A} سے \lr{5} سینٹی میٹر کے فاصلے پر نقاط کا لوسائی بنائیں۔}
\uq{362}{\lr{$m\overline{CD}=5.7\ cm$} کے قطعہ خط سے \lr{2.2} سینٹی میٹر کے فاصلے پر لوسائی بنائیں۔}
\uq{363}{\lr{$ABC=105^{o}$} کا زاویہ بنائیں۔ نقطہ \lr{P} کا لوسائی بنائیں جو \lr{$\overline{BA}$} اور \lr{$\overline{BC}$} سے برابر فاصلے پر حرکت کرتا ہو۔}
\uq{364}{دو نقاط \lr{E} اور \lr{F} کے درمیان فاصلہ \lr{5.4} سینٹی میٹر ہے۔ نقطہ \lr{P} کا لوسائی بنائیں جو \lr{E} اور \lr{F} سے برابر فاصلے پر حرکت کرتا ہے۔}
\uq{365}{ایک جزیرے پر دو شہر \lr{A} اور \lr{B} ہیں جن کا درمیانی فاصلہ \lr{8} کلو میٹر ہے۔ کاشف شہر \lr{A} سے \lr{6.8} کلو میٹر کے فاصلے پر جزیرے پر رہتا ہے اور شہر \lr{B} سے \lr{7.3} کلو میٹر کے فاصلے پر رہتا ہے۔ جزیرے پر ان نقاط کی نشاندہی کریں۔ جہاں کاشف وہ سکتا ہو۔}
\uq{366}{ایک مثلث \lr{CDE} بنائیں۔ جس میں \lr{$m\overline{CD}=7.6\ cm,\ m\angle D=45^{\circ}$} اور \lr{$m\overline{DE}=5.9\ cm$} ہو۔ تمام نقاط کا لوسائی کھینچیں جو کہ \lr{$\overline{CD}$} اور \lr{$\overline{CE}$} سے برابر فاصلے پر ہوں۔}
\uq{367}{مثلث \lr{LMN} بنائیں جس میں \lr{$m\overline{LM}=7\ cm,\ m\angle L=70^{\circ}$} اور \lr{$m\angle M=45^{\circ}$} ہو مثلث \lr{LMN} کے اندر ایک ایسا نقطہ معلوم کریں جو \lr{L} اور \lr{M} سے برابر فاصلے پر ہو اور \lr{L} سے \lr{3} سینٹی میٹر کے فاصلے پر ہو۔}
\uq{368}{ایک قائمۃ الزاویہ مثلث \lr{RST} بنائیں۔ جس میں \lr{$m\overline{RS}=6.8\ cm,\ m\angle S=90^{\circ}$} اور \lr{$m\overline{ST}=7.5\ cm$} ہو۔ مثلث \lr{RST} کے اندر ایک ایسا نقطہ معلوم کریں جو \lr{$\overline{RS}$} اور \lr{$\overline{RT}$} سے برابر فاصلے پر ہوں اور \lr{R} سے \lr{4.5} سینٹی میٹر کے فاصلے پر ہو۔}
\uq{369}{ایک مستطیل \lr{CVWX} بنائیں جس میں \lr{$m\overline{UV}=7.2\ cm$} اور \lr{$m\overline{VW}=5.6cm$} ہو۔ نقاط کا لوسائی بنائیں جو \lr{$\overline{UV}$} سے \lr{2} سینٹی میٹر اور \lr{W} سے \lr{3.5} سینٹی میٹر کے فاصلے پر ہو۔}
\uq{370}{تصور کریں دو سیل ٹاور نقاط \lr{A} اور \lr{B} محدداتی مستوی پر واقع ہیں۔ ایک آلہ \lr{GPS} مستوی پر کسی جگہ نصب ہے اور دونوں ٹاورز سے سگنل وصول کر رہا ہے۔ جہاز رانی کو یقینی بنانے کے لیے آلہ کو اس جگہ پر نصب کیا گیا ہے جو دونوں ٹاورز سے برابر فاصلے پر ہو۔ اس راستہ کا لوسائی بنائیں۔}
\uq{371}{ماہر وبائی امراض لوسائی کا استعمال انفیکشن زدہ علاقوں کی نشاندہی کرنے کے لیے کرتے ہیں۔ خاص طور پر متعدی بیماریوں کے لیے تاکہ پھیلاؤ کی پیشن گوئی اور حفاظتی اقدامات کیے جاسکیں۔ ایک بیماری کے پھیلاؤ کی صورت میں حکام ایک قرنطینہ علاقہ تعین کریں گے جو انفیکشن ذرائع سے \lr{10} کلو میٹر کے اندر ہو گا۔ بیماری کے پھیلاؤ کی نگرانی اور کنٹرول کے لیے قرنطینہ کے علاقے کی وضاحت کرتے ہوئے ماخذ سے \lr{10} کلو میٹر پر دو نقاط کا لوکس بنائیں۔}
\uq{372}{خزانہ ایک جزیرے پر کسی جگہ دفن ہے۔ خزانہ نقطہ \lr{A} سے \lr{24} کلو میٹر کے فاصلے پر ہے۔ اور نقاط \lr{B} اور \lr{C} سے ہم فاصلہ ہیں۔ پیمانہ \lr{10km=1cm} استعمال کرتے ہوئے وہ مقام معلوم کریں جہاں خزانہ دفن ہو سکتا ہے۔}
\uq{373}{سارہ کے گھر کے باغ میں ایک سیب کا درخت ہے۔ جو کیلے کے درخت سے \lr{90} میٹر فاصلے پر ہے۔ سارہ ایک آم کا درخت \lr{m} لگانا چاہتی ہے۔ جو سیب کے درخت سے \lr{64} میٹر اور کیلے کے درخت سے \lr{54} اور \lr{82} میٹر کے درمیانی فاصلے پر ہو۔ پیمانہ \lr{10m=1cm} استعمال کرتے ہوئے وہ مقام معلوم کریں جہاں آم کا درخت لگایا جائے۔}
\uq{374}{ایک مستطیل \lr{ABCD} بنائیں۔ جس میں \lr{$m\overline{AB}=5\ cm$} اور \lr{$m\overline{BC}=3.2\ cm$} ہو تمام نقاط کا لوسائی کھنچیں جو کہ نقطہ \lr{A} سے \lr{3.1} سینٹی میٹر کے فاصلے پر ہو۔ مستطیل کے اندر نقطہ \lr{P} پر ایک نشان لگائیں جو نقطہ \lr{A} سے \lr{3.1} سینٹی میٹر کے فاصلے پر ہو اور \lr{A} اور \lr{B} سے برابر فاصلے پر ہو}
\uq{375}{ایک مستطیل \lr{ABCD} بنائیں۔ جس میں \lr{$m\overline{AB}=5\ cm$} اور \lr{$m\overline{BC}=3.2\ cm$} ہو تمام نقاط کا لوسائی کھنچیں جو کہ نقطہ \lr{A} سے \lr{3.1} سینٹی میٹر کے فاصلے پر ہو۔ مستطیل کے اندر نقطہ \lr{P} پر ایک نشان لگائیں جو نقطہ \lr{A} سے \lr{3.1} سینٹی میٹر کے فاصلے پر ہو اور \lr{A} اور \lr{B} سے برابر فاصلے پر ہو}
\uq{376}{ایک متساوی الساقین مثلث \lr{DEF} بنائیں جس کا عمودی زاویہ نقطہ \lr{E} پر \lr{80$^{o}$} ہو۔ اور \lr{$m\overline{EF}=m\overline{DE}=4.8\ cm$} ہو۔ تمام نقاط کا لوسائی کھینچیں جو کہ: نقطہ \lr{E} سے \lr{2.8} سینٹی میٹر کے فاصلے پر ہو۔ مثلث کے اندر نقطہ \lr{X} کی نشاندہی کریں جو نقطہ \lr{E} سے \lr{2.8} سینٹی میٹر ہو اور \lr{$\overline{DE}$} اور \lr{$\overline{EF}$} مساوی فاصلے پر ہوں۔}
\uq{377}{ایک متساوی الساقین مثلث \lr{DEF} بنائیں جس کا عمودی زاویہ نقطہ \lr{E} پر \lr{80$^{o}$} ہو۔ اور \lr{$m\overline{EF}=m\overline{DE}=4.8\ cm$} ہو۔ تمام نقاط کا لوسائی کھنچیں جو کہ۔ \lr{$\overline{DE}$} اور \lr{$\overline{EF}$} سے مساوی فاصلے پر ہوں۔ مثلث کے اندر نقطہ \lr{X} کی نشان دہی کریں جو نقطہ \lr{E} سے \lr{2.8} سینٹی میٹر ہو اور \lr{$\overline{DE}$} اور \lr{$\overline{EF}$} مساوی فاصلے پر ہوں۔}
\uq{378}{ایک قائمۃ الزاویہ مثلث بنائیں جس کے اضلاع کی لمبایاں \lr{6 cm, 8 cm} اور \lr{10 cm} ہو۔}
\uq{379}{ایک مثلث \lr{ABC} بنائیں۔ جس میں \lr{$m\overline{AB}=5.3\ cm,\ m\angle A=30^{\circ}$} اور \lr{$m\angle B=120^{\circ}$} ہو۔ تمام نقاط کا لوسائی بنائیں جو \lr{A} اور \lr{B} سے برابر فاصلے پر ہوں۔}
\uq{380}{ایک مثلث بنائیں جس میں \lr{$m\overline{DE}=7.3\ cm,\ m\angle D=42^{\circ}$} اور \lr{$m\overline{EF}=5.4\ cm$} ہو۔}
\uq{381}{ایک مثلث \lr{XYZ} بنائیں جس میں \lr{$m\overline{YX}=8\ cm,\ m\overline{YZ}=7\ cm$} اور \lr{$m\overline{XZ}=6.5\ cm$} ہو۔ ان تمام نقاط کے لوسائی بنائیں جو \lr{$\overline{XY}$} اور \lr{$\overline{XZ}$} سے مساوی فاصلہ پر ہوں۔}
\uq{382}{ایک مثلث \lr{FGH} بنائیں جس میں \lr{$m\overline{FG}=m\overline{GH}=6.4\ cm,\ m\angle G=122^{\circ}$} ہو۔ تمام نقاط کو لوسائی بنائیں جو کہ۔ \lr{F} اور \lr{G} سے مساوی فاصلہ پر ہوں۔}
\uq{383}{ایک مثلث \lr{FGH} بنائیں جس میں \lr{$m\overline{FG}=m\overline{GH}=6.4\ cm,\ m\angle G=122^{\circ}$} ہو۔ تمام نقاط کو لوسائی بنائیں جو کہ \lr{$\overline{FG}$} اور \lr{$\overline{GH}$} سے مساوی فاصلہ پر ہوں۔}
\uq{384}{ایک مثلث \lr{FGH} بنائیں جس میں \lr{$m\overline{FG}=m\overline{GH}=6.4\ cm,\ m\angle G=122^{\circ}$} ہو۔ تمام نقاط کو لوسائی بنائیں جو کہ نقطہ کی نشان دہی کریں جہاں دونوں لوکس ایک دوسرے کو قطع کرتے ہیں}
\uq{385}{دو گھروں \lr{Q} اور \lr{R} کا درمیانی فاصلہ \lr{73} میٹر ہے۔ پیمانہ \lr{10m=1cm} استعمال کرتے ہوئے، نقطہ \lr{P} کا لوسائی بنائیں جو اس طرح حرکت کرتا ہے کہ \lr{Q} سے \lr{32} میٹر کے فاصلے پر ہو۔}
\uq{386}{دو گھروں \lr{Q} اور \lr{R} کا درمیانی فاصلہ \lr{73} میٹر ہے۔ پیمانہ \lr{10m=1cm} استعمال کرتے ہوئے، نقطہ \lr{P} کا لوسائی بنائیں جو اس طرح حرکت کرتا ہے کہ \lr{Q} اور \lr{R} کو ملانے والے خط سے \lr{48} میٹر کے فاصلے پر ہوں۔}
\uq{387}{ایک میدان مستطیل \lr{ABCD} کی شکل میں ہے جس میں \lr{$m\overline{AB}=70m$} اور \lr{$m\overline{BC}=60\ m$} ہو۔ پیمانہ \lr{10m=1cm} کا استعمال کرتے ہوئے مستطیل \lr{ABCD} بنائیں۔ میدان کے اندر وہ خطہ دکھائیں جو نقطہ \lr{C} سے \lr{30} میٹر سے کم فاصلے پر ہو اور \lr{$\overline{AB}$} سے \lr{25} میٹر سے زیادہ فاصلے پر ہو۔}
\uq{388}{نیچے دیے گئے تعددی تقسیم کے جدول سے کالمی نقشہ پر تعددی کثیر الاضلاع بنائیں:}
\dtab{|c|c|c|c|c|c|}{30-34 & 25-29 & 20-24 & 15-19 & 10-14 & \textbf{Weight(kg)}\\ 22 & 30 & 23 & 17 & 06 & \textbf{Frequency(f)}}
\uq{389}{دسویں جماعت کے طلبہ کے ریاضی میں حاصل کردہ نمبروں کی تعددی تقسیم درج ذیل جدول میں دی گئی ہے:}
\dtab{c|c|c|c|c|c|c}{68-72 & 62-67 & 55-61 & 45-54 & 38-44 & 35-37 & \urdu{\textbf{نمبرز}}\\ 3 & 9 & 13 & 16 & 2 & 1 & \urdu{\textbf{تعدد}}}
\usc{اوپر دیئے گئے تعددی تقسیم کے جدول کو بذریعہ کالمی نقشہ ظاہر کریں۔}
\uq{390}{نیچے دیے گئے گروہی مواد سے کالمی نقشہ پر تعددی کثیر الاضلاع بنائیں:}
\dtab{|c|c|c|c|c|c|}{25-27 & 20-25 & 12-20 & 8-12 & 5-8 & \urdu{\textbf{جماعتی حدود}}\\ 14 & 32 & 25 & 12 & 2 & \urdu{\textbf{تعدد}} \lr{\textbf{(f)}}}
\uq{391}{سکول کے سٹاف کے لیے کرسیوں کی مرمت پر درج ذیل اخراجات (سینکڑوں روپوں میں) کی ضرورت ہے؟}
\nums{145\quad152\quad153\quad156\quad158\quad160\quad146\quad152\quad155\quad159\\ 161\quad163\quad165\quad147\quad148\quad151\quad154\quad156\quad158\quad160\\ 144\quad167\quad151\quad150\quad152\quad149\quad145\quad153\quad152\quad155}
\usc{جماعتوں کی تعداد \lr{8} لے کر ٹیلی نشان کے طریقے سے ایک تعددی تقسیم کا جدول بنائیں اور آخری تین جماعتوں کے تعدادت لکھیں۔}
\uq{392}{ذیل میں ایک ہائی سکول کے \lr{30} طلبہ کے اوزان کلو گرام میں دیے گئے ہیں:}
\nums{30\quad33\quad24\quad21\quad15\quad39\quad37\quad44\quad42\\ 33\quad28\quad29\quad32\quad31\quad28\quad26\quad32\quad34\\ 38\quad36\quad41\quad30\quad35\quad41\quad23\quad26\quad18}
\usc{جماعتوں کی تعداد \lr{5} لے کر ایک تعددی تقسیم کا جدول بنائیں اور ایک کثیر الاضلاع بنائیں۔}
\uq{393}{دسویں جماعت کے طلبہ کے ایک گروپ نے انگریزی کے مضمون میں \lr{(100} میں سے) درج ذیل نمبرز حاصل کیے:}
\nums{58\quad59\quad58\quad33\quad40\quad58\quad45\quad46\quad43\quad45\quad45\\ 50\quad52\quad49\quad50\quad57\quad52\quad55\quad49\quad50\quad62\quad49\\ 48\quad44\quad42\quad47\quad46\quad47\quad46\quad53\quad40\quad44}
\usc{جماعتوں کی تعداد \lr{6} لے کر مواد کی ایک تعددی تقسیم کا جدول بنائیں۔ نیز سب سے کم تعدد والی جماعتی حد معلوم کریں اور اس مواد کا کالمی نقشہ بھی بنائیں۔}
\uq{394}{ایک ہائی سکول کے \lr{30} اساتذہ کو ایک ہفتہ میں جو ٹیلی فون کالز موصول ہوئیں وہ درج ذیل ہیں:}
\nums{5\quad8\quad11\quad25\quad13\quad16\quad20\quad17\quad15\quad16\quad30\quad21\quad14\quad8\\ 6\quad22\quad26\quad15\quad19\quad35\quad29\quad31\quad23\quad25\quad20\quad10\quad9\quad7}
\usc{جماعتوں کی تعداد \lr{7} لے کر تعددی تقسیم بنائیں۔}
\uq{395}{ایک علاقے کے \lr{76} افراد کی عمر (سالوں میں) کی تعددی تقسیم دی گئی ہے۔ اس مواد کے لیے کالمی نقشہ بنائیں:}
\dtab{|c|c|c|c|c|c|c|c|}{27-30 & 20-27 & 17-20 & 12-17 & 9-12 & 4-9 & 2-4 & \urdu{\textbf{جماعتی حدود}}\\ 4 & 7 & 10 & 20 & 18 & 10 & 7 & \urdu{\textbf{تعدد}} \lr{\textbf{(f)}}}
\uq{396}{\lr{30} طلبہ نے ریاضی کے مضمون میں اپنے سالانہ امتحان میں \lr{(100} میں سے) درج ذیل نمبر حاصل کیے۔ نیچے دی گئی تعددی تقسیم کے لیے تعددی کثیر الاضلاع بنائیں:}
\dtab{|c|c|c|c|c|c|}{56-69 & 60-64 & 55-59 & 50-54 & 45-49 & \urdu{\textbf{نمبرز}}\\ 3 & 7 & 4 & 9 & 2 & \urdu{\textbf{تعدد}}}
\uq{397}{درج ذیل انگریزی مضمون میں طلبہ کے حاصل کردہ نمبر دیے گئے ہیں}
\dtab{|c|c|c|c|c|c|}{35-39 & 30-34 & 25-29 & 20-24 & 15-19 & \urdu{حاصل کردہ نمبر}\\ 5 & 17 & 35 & 18 & 9 & \urdu{تعدد}}
\usc{ان نمبروں کا حسابی وسط (براہ راست اور مختصر طریقے سے) معلوم کریں۔}
\uq{398}{درج ذیل انگریزی مضمون میں طلبہ کے حاصل کردہ نمبر دیے گئے ہیں}
\dtab{|c|c|c|c|c|}{30-34 & 25-29 & 20-24 & 15-19 & \urdu{حاصل کردہ نمبر}\\ 17 & 35 & 18 & 9 & \urdu{تعدد}}
\usc{ان نمبروں کا وسطانیہ معلوم کریں۔}
\uq{399}{درج ذیل تعددی تقسیم دی گئی ہے اس تعددی تقسیم کا عادہ معلوم کریں:}
\dtab{|c|c|c|c|c|c|}{25-29 & 20-24 & 15-19 & 10-14 & 5-9 & \urdu{\textbf{جماعتی حدود}}\\ 2 & 11 & 18 & 8 & 1 & \urdu{\textbf{تعدد}}}
\uq{400}{ایک ہوسٹل میں \lr{10} طلبہ کے دوپہر کے کھانے کی ماہانہ حاضریوں کا ریکارڈ کچھ اس طرح ہے: \lr{21,15,16,18,14,17,15,12,13,11}۔ حاضریوں کا وسطانیہ اور عادہ معلوم کریں اگر \lr{D= A-20} فرض کریں تو حسابی اوسط بھی معلوم کریں۔}
\uq{401}{تقسیم انعامات کے دن \lr{50} طلبہ کے لائے ہوئے جیب خرچ کی تفصیل درج ذیل جدول میں دی گئی ہے:}
\dtab{|c|c|c|c|c|c|}{25-30 & 20-25 & 15-20 & 10-15 & 5-10 & \urdu{روپے}\\ 4 & 7 & 18 & 9 & 12 & \urdu{تعدد} \lr{(f)}}
\usc{دیے گئے مواد کا حسابی وسط مختصر طریقے سے معلوم کریں۔}
\uq{402}{سکول میں اساتذہ کے ایک گروپ نے کرکٹ میچ منعقد کیا جس میں تین بچوں حارث، ماہم اور منال نے درج ذیل دوڑیں بنائیں:}
\dtab{|c|c|c|c|c|}{85 & 70 & 55 & 50 & \urdu{\textbf{حارث کی دوڑیں}}\\ 45 & 60 & 60 & 75 & \urdu{\textbf{ماہم کی دوڑیں}}\\ 42 & 66 & 77 & 80 & \urdu{\textbf{منال کی دوڑیں}}}
\usc{یہ فیصلہ کیا گیا کہ جو بچہ سب سے زیادہ اوسط دوڑیں بنائے گا اسے \lr{1000} روپے انعام دیے جائیں گے۔ انعام کی رقم کون حاصل کرے گا؟}
\uq{403}{ایک معلومات عامہ کے پروگرام میں دو ساتھیوں حفصہ اور فاطمہ نے حصہ لیا اور انہوں نے درج ذیل پوائنٹس حاصل کیے۔ \lr{46,48,74,61,58,51,45} اور \lr{50}۔ اوسط پوائنٹس بذریعہ انحراف \lr{D=x=-58} معلوم کریں۔}
\uq{404}{ایک شخص نے درج ذیل کھانے کی اشیا خریدیں:}
\dtab{|c|c|c|}{\urdu{\textbf{لاگت (روپوں میں)}} & \urdu{\textbf{مقدار (کلوگرام میں)}} & \urdu{\textbf{کھانے کی اشیا}}\\ \hline 96 & 10 & \urdu{چاول}\\ 48 & 12 & \urdu{آٹا}\\ 190 & 4 & \urdu{گھی}\\ 49 & 3 & \urdu{چینی}\\ 650 & 2 & \urdu{گوشت}}
\usc{کھانے والی اشیا کا وزنی وسطی کلو گرام کیا ہے؟}
\uq{405}{\lr{20} لڑکوں کی عمروں کا حسابی اوسط \lr{13} سال \lr{4} ماہ اور \lr{5} دن ہے۔ ان کی عمروں کا مجموعہ معلوم کریں۔ اگر ان میں سے ہر ایک لڑکے کی عمر \lr{15} سال ہو تو باقی لڑکوں کی اوسط عمر کیا ہے؟}
\uq{406}{نیچے دیا گیا ایک تعددی تقسیم کا جدول ہے جسے انحرافات \lr{D=X-20} سے بنایا گیا ہے۔ اس مواد سے حسابی اوسط معلوم کریں:}
\dtab{c|c|c|c|c|c|c}{4 & 2 & 0 & -2 & -4 & -6 & \textbf{D}\\ 12 & 26 & 20 & 6 & 3 & 1 & \textbf{f}}
\uq{407}{جمال کے ریاضی میں \lr{8} ماہانہ ٹیسٹوں میں حاصل کردہ نمبر \lr{75,76,80,80,82,82,82,85} تھے۔ حاصل کردہ نمبروں کا عادہ معلوم کریں۔}
\uq{408}{\lr{10} طلبہ سے پوچھا گیا کہ انہوں نے پچھلے ہفتے کل \lr{20} سوالات میں سے کتنے سوالات حل کیے۔ ریکارڈ کیے گئے سوالات کے حل کی تعداد \lr{13,14,15,11,16,10,19,20,18,17} تھی۔ مواد کا عادہ معلوم کریں۔}
\uq{409}{ایک سکول کے \lr{15} طلبہ سے ایک سروے کیا گیا اور طلبہ سے ان کے پسندیدہ رنگ کے بارے میں پوچھا گیا۔ جوابات یہ ہیں۔ جامنی، پیلا، جامنی، پیلا، پیلا، سرخ، نیلا، سبز، پیلا، پیلا، سرخ، نیلا، پیلا، جامنی، سبز۔ مواد کا عادہ معلوم کریں۔}
\uq{410}{ایک اچھی شہرت کے حامل سکول سے \lr{10} طلبہ پانچویں جماعت کے فریق \lr{A} اور \lr{10} طلبہ فریق \lr{B} سے لیے گئے۔ ان کہ اوزان کلو گرام میں پیمائش کیے گئے اور انہیں درج ذیل تفصیل میں ریکارڈ کیا گیا:}
\dtab{|c|c|c|c|c|c|c|c|c|c|}{40 & 33 & 31 & 34 & 35 & 29.5 & 32 & 28 & 30 & \textbf{Weights(kg)SectionA}\\ 6.5 & 36 & 39.5 & 38 & 32.8 & 35 & 34.5 & 31.5 & 35 & \textbf{Weights(kg)SectionB}}
\usc{فریق \lr{A} اور فریق \lr{B} کے لیے وسطاً اوزان معلوم کریں۔}
\uq{411}{ایک حکومت ایک سکول کے نیچے دیے گئے پانچ شعبوں کے لیے \lr{200,000} روپے کے فنڈز مختص کرتی ہے۔ \lr{(i)} سکول لائبریری کے لیے \lr{35,000} روپے، \lr{(ii)} کھیل کی سہولیات: \lr{25,000} روپے، \lr{(iii)} پارکنگ کی جگہ \lr{40,000} روپے، \lr{(iv)} کمروں کی مرمت \lr{45,000} روپے، \lr{(v)} فرنیچر \lr{55,000} روپے۔ سکول کے ہر شعبے میں فنڈز مختص کرنے کی اوسط معلوم کریں۔}
\uq{412}{ایک امتحان میں \lr{100} طلبہ کے حاصل کردہ نمبرز \lr{(100} میں سے) دیے گئے ہیں۔ ان کے نمبروں کا حسابی اوسط معلوم کریں:}
\dtab{|c|c|c|c|c|c|}{50-55 & 45-50 & 40-45 & 35-40 & 30-35 & \urdu{\textbf{نمبرز}}\\ 18 & 23 & 18 & 16 & 14 & \urdu{\textbf{طلبہ کی تعداد}}}
\uq{413}{درج ذیل میں \lr{12.5} میں سے \lr{10} مختلف انحرافات دیے گئے ہیں۔ حسابی اوسط معلوم کریں۔ \lr{6,-4.2,-6.8,5.5,14,11.3,-2,3.5,9,8.7}}
\uq{414}{درج ذیل تعددی تقسیم میں سکول کے \lr{200} لڑکوں کے قد (انچوں میں) ریکارڈ کیے گئے ہیں۔ مختصر طریقے سے لڑکوں کا اوسط قد معلوم کریں:}
\dtab{|c|c|c|c|c|c|c|c|c|c|}{59 & 58 & 57 & 56 & 55 & 54 & 53 & 52 & 51 & \urdu{\textbf{قد (x) (انچوں میں)}}\\ 6 & 16 & 38 & 45 & 55 & 24 & 8 & 5 & 2 & \urdu{\textbf{تعدد (f)}}}
\uq{415}{\lr{100} کھلاڑیوں کے قد تقریباً (انچوں میں) درج ذیل جدول میں دیے گئے ہیں۔ وسطانیہ معلوم کریں۔}
{\scriptsize\setlength{\tabcolsep}{3pt}\dtab{|c|c|c|c|c|c|c|c|c|c|}{70.5-71.5 & 69.5-70.5 & 68.5-69.5 & 67.5-68.5 & 66.5-67.5 & 65.5-66.5 & 64.5-65.5 & 63.5-64.5 & 62.5-63.5 & \urdu{\textbf{قد (انچوں میں)}}\\ 2 & 3 & 12 & 13 & 30 & 20 & 10 & 6 & 4 & \urdu{\textbf{طلبہ کی تعداد}}}}
\uq{416}{درج ذیل میں \lr{50} آدمیوں کے اوزان (کلوگرام) میں دیے گئے ہیں۔ اوزان کا وسطانیہ معلوم کریں:}
\dtab{|c|c|c|c|c|c|}{130-134 & 125-129 & 120-124 & 115-119 & 110-114 & \urdu{\textbf{اوزان (کلوگرام میں)}}\\  & 6 & 23 & 12 & 5 & \urdu{\textbf{آدمیوں کی تعداد (f)}}}
\uq{417}{درج ذیل آٹھویں جماعت کے \lr{40} طلبہ کے قد (انچوں میں) دیے گئے ہیں۔ مواد کا عادہ معلوم کریں۔}
\dtab{|c|c|c|c|c|c|c|}{58-60 & 56-58 & 54-56 & 52-54 & 50-52 & 48-50 & \urdu{\textbf{قد (انچوں میں)}}\\ 3 & 6 & 9 & 10 & 7 & 5 & \urdu{\textbf{طلبہ کی تعداد (f)}}}
\uq{418}{ایک دوا ساز کمپنی نے شہر کے سات مختلف علاقوں میں دوائی کے نمونے کی مارکیٹنگ شروع کی۔ کمپنی نے دوا کے پیکٹ شہر کے ہر علاقہ میں تقسیم کیے اور ہر علاقے کی مانگ کے مطابق دوائی تقسیم کی۔ دیے گئے مواد کا حسابی اوسط اور اوزانی اوسط معلوم کریں۔}
\dtab{|c|c|c|}{\urdu{\textbf{اوزان (کلوگرام میں)}} & \urdu{\textbf{پیکٹوں کی تعداد}} & \urdu{\textbf{شہر کے علاقے}}\\ \hline 5 & 15 & A\\ 4 & 25 & B\\ 3 & 18 & C\\ 4 & 23 & D\\ 2 & 15 & E\\ 1 & 10 & F\\ 2 & 8 & G}
\uq{419}{ایک کھلونا بنانے والی فیکٹری نے ایک مہینے میں جتنے کھلونے فروخت کیے، درج ذیل مواد پر غور کریں۔}
\dtab{|c|c|c|c|c|c|}{50-60 & 40-50 & 30-40 & 20-30 & 10-20 & \urdu{\textbf{جماعتی حدود}}\\ 20 & 29 & 45 & 28 & 15 & f}
\usc{(الف) فیکٹری سے فروخت کیے گئے کھلونوں کی تعداد کا حسابی اوسط وسطانیہ اور عادہ معلوم کریں۔}
\usc{(ب) نیز اس جدول کی عادہ جماعت بھی بتائیں۔}
\uq{420}{درج ذیل میں \lr{40} طلبہ کے اوزان تقریباً (پونڈ میں) ریکارڈ کیے گئے ہیں۔}
\nums{138\quad164\quad150\quad132\quad144\quad125\quad149\quad157\quad146\\ 140\quad147\quad136\quad148\quad152\quad144\quad168\quad126\quad138\\ 163\quad119\quad154\quad165\quad146\quad173\quad142\quad147\quad135\\ 140\quad135\quad161\quad145\quad135\quad142\quad150\quad156\quad145}
\usc{جماعتوں کی تعداد \lr{6} لے کر ایک تعددی تقسیم کا جدول بنائیں۔ (ب) کالمی نقشہ بنائیں۔ (ج) دیے ہوئے مواد کا تعددی کثیر الاضلاع بھی بنائیں۔}
\uq{421}{ملحقہ تقسیم میں زیادہ سے زیادہ بوجھ کلو گرام میں دکھایا گیا ہے۔ جسے کچھ رسیوں کی مدد سے سہارا دیا گیا ہے۔ مختصر طریقہ کار استعمال کرتے ہوئے اوسط وزن معلوم کریں:}
\dtab{|c|c|c|c|c|c|c|}{118-122 & 113-117 & 108-112 & 103-107 & 98-102 & 93-97 & \urdu{\textbf{زیادہ سے زیادہ بوجھ (کلوگرام میں)}}\\ 2 & 6 & 12 & 8 & 5 & 2 & \urdu{\textbf{رسیوں کی تعداد}}}
\uq{422}{دیے گئے تعددی تقسیم کے جدول کا کالمی نقشہ اور تعددی کثیر الاضلاع بنائیں:}
\dtab{|c|c|c|c|c|c|}{25-29 & 20-24 & 15-19 & 10-14 & 5-9 & \urdu{\textbf{جماعتی حدود}}\\ 2 & 11 & 18 & 8 & 1 & \urdu{\textbf{تعدد}}}
\uq{423}{دیے گئے جدول کا اوزانی اوسط معلوم کریں:}
\dtab{|c|c|c|}{\urdu{\textbf{اشیا کی لاگت (روپوں میں)}} & \urdu{\textbf{مقداریں}} & \urdu{\textbf{اشیا}}\\ \hline 500 & 20 & \urdu{کرسی}\\ 400 & 20 & \urdu{میز}\\ 750 & 10 & \urdu{تختہ سیاہ}\\ 230 & 25 & \urdu{ٹیوب لائٹ}\\ 950 & 09 & \urdu{الماری}}
\uq{424}{دو شراکت دار اسلم اور کلثوم ایک کمپنی چلاتے ہیں۔ درج ذیل مواد میں کمپنی میں کام کرنے والی ملازمین کی ہفتہ وار اجرت (روپوں میں) دی گئی ہے۔}
\dtab{|c|c|c|c|c|c|}{1100- & 900-1000 & 800-900 & 700-800 & 600-700 & \urdu{\textbf{اجرت (روپوں میں)}}\\ 1 & 21 & 7 & 5 & 3 & \urdu{\textbf{ملازمین}}}
\usc{دیے گئے جدول کا حسابی وسط وسطانیہ اور عادہ معلوم کریں۔}
\uq{425}{نیچے دیے گئے تعددی تقسیم کے جدول کی مدد سے کالمی نقشہ پر تعددی کثیر الاضلاع بنائیں:}
\dtab{|c|c|c|c|c|c|c|}{76-80 & 73-75 & 65-72 & 60-64 & 57-59 & 50-56 & \urdu{\textbf{وزن (کلوگرام میں)}}\\ 8 & 15 & 30 & 40 & 32 & 25 & \urdu{\textbf{تعدد (f)}}}
\uq{426}{ایک سکول کا پرنسپل پانچ مختلف شعبوں کے لیے \lr{50,000} روپے کے فنڈز مختص کرتا ہے۔ \lr{(i)} کرسیاں \lr{= 15000} روپے \lr{(ii)} میز \lr{= 12000} روپے \lr{(iii)} بلیک بورڈ \lr{= 6000} روپے \lr{(iv)} کمروں کی تزئین و آرائش \lr{= 10000} روپے \lr{(v)} باغبانی \lr{= 7000} روپے۔ سکول کے ہر شعبے میں فنڈز مختص کرنے کی اوسط معلوم کریں۔}
\uq{427}{اسلم ایک ڈائس پھینکتا ہے۔ ہندسہ \lr{3} یا \lr{4} حاصل ہونے کا احتمال کیا ہو گا؟ نیز ہندسہ \lr{3} یا \lr{4} حاصل نا ہونے کا احتمال بھی معلوم کریں۔}
\uq{428}{عبد الہادی نے کارڈز کو \lr{1} تا \lr{30} تک لیبل کیا اور ایک ڈبہ میں ڈال دیا۔ وہ ایک کارڈ بے ترتیب منتخب کرتا ہے۔ منتخب شدہ کارڈ کا احتمال کیا ہے جب کہ وہ:}
\usc{\lr{(i)} عدد \lr{25} ہو \lr{(ii)} عدد \lr{17} سے \lr{22} کے درمیان ہو \lr{(iii)} عدد کم از کم \lr{20} ہو}
\usc{\lr{(iv)} عدد \lr{27} اور \lr{29} نا ہو \lr{(v)} عدد \lr{12} سے \lr{15} کے درمیان نا ہو۔}
\uq{429}{عائشہ کے امتحان پاس کرنے کا احتمال \lr{0.85} ہے۔ عائشہ کے امتحان پاس نا کرنے کا احتمال کیا ہو گا؟}
\uq{430}{اچھی طرح ملائے گئے تاش کے \lr{52} پتوں میں سے ایک پتہ منتخب کیا گیا۔ احتمال کیا ہو گا اگر منتخب شدہ پتہ: \lr{(i)} ملکہ ہو \lr{(ii)} نا ملکہ اور نا ہی غلام ہو۔}
\uq{431}{تاش کے \lr{52} پتوں میں سے ایک پتہ منتخب کیا گیا۔ احتمال کیا ہو گا اگر منتخب شدہ پتہ: \lr{(i)} غلام ہو \lr{(ii)} اینٹ نا ہو۔}
\uq{432}{تایش ایک سکہ اچھالتا ہے اور ایک ڈائس پھینکتا ہے۔ درج ذیل واقعات کا احتمال کیا ہو گا؟}
\usc{\lr{(i)} سکہ پر تیل اور ڈائس پر کم از کم ہندسہ \lr{4} آئے۔}
\usc{\lr{(ii)} سکہ پر ہیڈ اور ڈائس پر ہندسے \lr{2} اور \lr{3} آئیں۔}
\usc{\lr{(iii)} سکہ پر ہیڈ اور ٹیل اور ڈائس پر ہندسہ \lr{6} آئے۔}
\usc{\lr{(iv)} سکہ پر ٹیل نا آئے اور ڈائس پر ہندسہ \lr{5} نا آئے۔}
\usc{\lr{(v)} سکہ پر ہیڈ اور ڈائس پر ہندسے \lr{5} اور \lr{2} نا آئیں۔}
\uq{433}{اگر \lr{3} میزائل \lr{A، B} اور \lr{C} ہوں اور انہیں ایک ٹارگٹ پر فائر کرنا ہو۔ اگر ٹارگٹ فائر کرنے کا احتمال \lr{$P(A)=\frac{1}{4},\ P(B)=\frac{3}{7},\ P(C)=\frac{5}{9}$} ہو تو احتمال معلوم کریں۔}
\usc{\lr{(1)} میزائل \lr{A} ٹارگٹ کو نشانہ نہیں بناتا}
\usc{\lr{(iii)} میزائل \lr{B} ٹارگٹ کو نشانہ نہیں بناتا}
\usc{\lr{(ii)} میزائل \lr{C} ٹارگٹ کو نشانہ نہیں بناتا}
\uq{434}{ایک تھیلے میں \lr{5} نیلی گیندیں اور \lr{8} سبز گیندیں ہیں۔ ایک گیند کے منتخب ہونے کا احتمال معلوم کریں جب وہ:}
\usc{\lr{(i)} نیلی گیند ہو \lr{(ii)} سبز گیند ہو \lr{(iii)} سبز گیند نہ ہو}
\uq{435}{ایک محقق نے روسی فوج کی ہارس ریکس سے ہونے والی ہلاکتوں کی تعداد کا مواد جمع کیا۔ جدول درج ذیل ہے۔}
\dtab{|c|c|c|c|c|c|c|c|}{6 & 5 & 4 & 3 & 2 & 1 & 0 & \urdu{\textbf{ہلاکتوں کی تعداد}}\\ 10 & 15 & 32 & 40 & 87 & 50 & 60 & \urdu{\textbf{تعداد}}}
\usc{دیے گئے مواد کی نسبتی تعداد معلوم کریں۔}
\uq{436}{\lr{750} نمونوں میں سے ناقص مصنوعات کا جدول درج ذیل دیا گیا ہے۔ گنے جدول کے لیے نسبتی تعداد معلوم کریں۔}
\dtab{|c|c|c|c|c|c|c|c|}{6 & 5 & 4 & 3 & 2 & 1 & 0 & \urdu{\textbf{ناقص مصنوعات کی تعداد فی نمونہ}}\\ 40 & 50 & 105 & 85 & 94 & 40 & 120 & \urdu{\textbf{نمونوں کی تعداد}}}
\uq{437}{واقفیت عامہ پر ایک کوئز کا مقابلہ منعقد کیا گیا۔ سوالات کے \lr{100} سیٹوں کے لیے \lr{5} سوالات میں سے درست جوابات کی تعداد ذیل میں دی گئی ہے۔}
\dtab{|c|c|c|c|c|c|c|}{5 & 4 & 3 & 2 & 1 & 0 & \textbf{x}\\ 9 & 18 & 25 & 15 & 23 & 10 & \textbf{f}}
\usc{دیے گئے مواد کے لیے نسبتی تعداد معلوم کریں۔}
\uq{438}{ایک جماعت کے \lr{108} طلبہ سے ایک سروے کیا گیا۔ اور ان کے پسندیدہ کھانے کے بارے میں پوچھا گیا۔ جوابات درج ذیل ہیں۔}
\dtab{|c|c|c|c|c|}{\urdu{بابی کیو} & \urdu{مرغی} & \urdu{تازہ جوس} & \urdu{بریانی} & \urdu{\textbf{کھانے کی اشیا کام}}\\ 15 & 21 & 07 & 40 & \urdu{\textbf{طلبہ کی تعداد}}}
\usc{\lr{(i)} کتنے فی صد طلبہ بریانی کو پسند کرتے ہیں؟}
\usc{\lr{(ii)} کتنے فی صد طلبہ مرغی پسند کرتے ہیں؟}
\usc{\lr{(iii)} طلبہ نے کون سا کھانا سب سے کم پسند کیا؟}
\usc{\lr{(iv)} طلبہ کون سے کھانے کو سب سے زیادہ ترجیح دیتے ہیں؟}
\uq{439}{اگر ایک تجربہ \lr{200} بار دہرایا جاتا ہے۔ تو دیے گئے مواد کی متوقع تعداد معلوم کریں۔}
\dtab{|c|c|c|c|c|c|c|c|}{6 & 5 & 4 & 3 & 2 & 1 & 0 & \textbf{x}\\ 0.07 & 0.17 & 0.09 & 0.18 & 0.17 & 0.21 & 0.11 & \textbf{P(x)}}
\uq{440}{\lr{6} ڈائس اچھالتے ہوئے \lr{5} مرتبہ چھکا آنے کا احتمال \lr{$\frac{2}{5}$} ہے۔ ڈائس \lr{200} پھینکے گئے ہیں۔ آپ منی بار توقع کریں گے کہ یہ \lr{5} چھکے دکھائے گا؟}
\uq{441}{دیے گئے جدول کا نسبتی تعداد معلوم کریں۔}
\dtab{|c|c|c|c|c|c|c|c|}{8 & 7 & 6 & 5 & 4 & 3 & 2 & \textbf{x}\\ 2 & 8 & 10 & 9 & 6 & 5 & 3 & \textbf{f}}
\uq{442}{عبد الرحمن نے \lr{100} نمبروں میں سے مختلف مضامین میں مختلف نمبر حاصل کیے۔ تفصیل درج ذیل ہے۔}
\dtab{|c|c|c|c|c|c|c|}{\urdu{کمپیوٹر سائنس} & \urdu{سائنس} & \urdu{ریاضی} & \urdu{اسلامیات} & \urdu{اردو} & \urdu{انگریزی} & \urdu{\textbf{مضمون}}\\ 85 & 81 & 95 & 72 & 80 & 75 & \urdu{\textbf{حاصل کردہ نمبرز}}}
\usc{مندرجہ بالا مواد کا نسبتی تعداد معلوم کریں۔}
\uq{443}{چھ ڈائس \lr{50} مرتبہ پھینکے گئے۔ مختلف تعداد میں \lr{6} آنے کے احتمال درج ذیل دیے گئے ہیں۔ درج ذیل مواد کی متوقع تعداد معلوم کریں۔}
\dtab{|c|c|c|c|c|c|c|c|}{6 & 5 & 4 & 3 & 2 & 1 & 0 & \textbf{x}\\ 0.15 & 0.20 & 0.10 & 0.24 & 0.12 & 0.10 & 0.09 & \textbf{P(x)}}
\usc{ہر چھکے کو وقوع پذیر ہونے کے متوقع تعداد معلوم کریں۔}
\uq{444}{ایک صراحی میں \lr{10} سرخ گیندیں، \lr{5} سبز گیندیں اور \lr{8} ٹیلی گیندیں ہیں۔ بے ترتیب چناؤ کرتے ہوئے احتمال معلوم کریں:}
\usc{\lr{(i)} سبز گیند ہو \lr{(ii)} سرخ گیند ہو \lr{(iii)} نیلی گیند ہو \lr{(iv)} سرخ گیند نا ہو \lr{(v)} سبز گیند نا ہو}
\uq{445}{تین سکے ایک ساتھ اچھالے گئے۔ احتمال معلوم کریں اگر:}
\usc{\lr{(i)} تین ہیڈ ہوں \lr{(ii)} کم از کم دو ٹیل ہوں \lr{(iii)} کم از کم دو ہیڈ نا ہوں \lr{(iv)} دو ہیڈ نا ہوں}
\uq{446}{تاش کے \lr{52} پتوں میں سے ایک پتہ منتخب کیا گیا۔ احتمال کیا ہو گا اگر:}
\usc{\lr{(i)} سرخ رنگ کا بادشاہ اور غلام ہو۔ \lr{(ii)} \lr{2} چڑیا اور حکم کا \lr{2} نا ہو۔}
\uq{447}{چھ سکے \lr{600} بار اچھالے گئے۔ نیچے دیے گئے جدول میں ٹیل کے وقوع پذیر ہونے کا ریکارڈ دیا گیا ہے۔}
\dtab{|c|c|c|c|c|c|c|c|}{6 & 5 & 4 & 3 & 2 & 1 & 0 & \urdu{\textbf{ٹیل کی تعداد}}\\ 16 & 123 & 76 & 80 & 105 & 90 & 110 & \urdu{\textbf{تعدد}}}
\usc{دیے گئے جدول کا نسبتی تعداد معلوم کریں۔}
\uq{448}{\lr{25} اشیا پر مشتمل لاٹ میں سے \lr{18} اشیا ناقص ہیں۔ غیر ناقص اشیا کا نسبتی تعداد معلوم کریں اور غیر ناقص اشیا کی متوقع تعداد بھی معلوم کریں۔}

\end{document}
